\chapter{Plutus Core on Cardano}
\label{chap:cardano-uplc}
\section{Protocol versions}
The Cardano blockchain controls the introduction of features through the use of
\emph{protocol versions}, a field in the protocol parameters.  The major
protocol version is used to indicate when forwards-incompatible changes
(i.e. those that allow blocks that were not previously allowed) are made to the
rules of the chain.  This is a hard fork of the chain.

The protocol version is part of the history of the chain, as are all protocol
parameters.  That means that all blocks are associated with the protocol version
from when they were created, so that they can be interpreted correctly.

In summary, conditioning on the protocol version is the main way in which we can introduce changes in behaviour.

Table~\ref{table:protocol-versions} lists the protocol versions that are
relevant to the use of Plutus Core on Cardano.  This table is based on the
\texttt{features-table.md} document attached to CIP-0059~\cite{CIP-0059}.  For
brevity, in the remainder of this document we will sometimes abbreviate
``Protocol Version'' to PV and omit the minor version number, so for example PV9
means Protocol Version 9.0.

The table includes an entry for Protocol Version 12.0, which will initiate
Cardano's Dijkstra era, and this document uses the name ``Dijkstra HF'' for the
corresponding hard fork.  The date of this hard fork is unknown at the time of
writing, and its actual name will probably be different.


\begin{table}[H]
  \centering
    \begin{tabular}{|l|l|r|r|l|}
        \hline
        \thead{Protocol version} & \thead{Hard fork} & \thead{Epoch} & \thead{Slot} & \thead{Date} \\
        \hline
        5.0  & Alonzo      & 290 &   39916975  & September 2021 \\
        7.0  & Vasil       & 365 &   72316896  & June 2022 \\
        8.0  & Valentine   & 394 &   84844885  & February 2023 \\
        9.0  & Chang       & 507 &  133660855  & September 2024 \\
        10.0 & Plomin      & 537 &  146620809  & January 2025 \\
        11.0 & van Rossem  & 644 &  192845302  & July 2026 \\
        12.0 & Dijkstra HF &     &             & TBD\\
        \hline
    \end{tabular}
    \caption{Protocol versions}
    \label{table:protocol-versions}
\end{table}


\section{Ledger languages}

The Cardano ledger uses Plutus Core as the programming language for
\emph{scripts}.  The ledger in fact supports multiple different interpretations
for scripts, and so each script is tagged with a \emph{ledger language} that
tells the ledger how to interpret it.  Since the ledger must always be able to
evaluate old scripts and get the same answer, the ledger language must pin down
everything about how the script is evaluated, including:

\begin{enumerate}
  \item How to interpret the script itself (e.g. as a Plutus Core program, what
    versions of the Plutus Core language are allowable)
  \item Other configuration the script may need in order to run (e.g. the set of
    builtin types and functions and their interpretations, cost model
    parameters)
  \item How the script is invoked (e.g. after having certain arguments passed to
    it)
\end{enumerate}

There are currently three ``Plutus'' ledger languages (ie, ledger languages
whose underlying programming language is Plutus Core) in use on
Cardano:\footnote{ Note that ledger languages are completely distinct from the
point of view of the ledger, the ``V1''/``V2'' naming is suggestive of the fact
that these two ledger languages are related, but in the implementation they are
completely independent.  },
\begin{enumerate}
  \item \LL{PlutusV1}
  \item \LL{PlutusV2}
  \item \LL{PlutusV3}
\end{enumerate}

\noindent Table~\ref{table:ll-introduction} shows when each Plutus ledger
language was introduced.  Ledger languages remain available permanently after
they have been introduced.  A fourth ledger language, \LL{PlutusV4}, will become
available at Protocol Version 12.0; this will mostly differ from \LL{PlutusV3}
in the API for interacting with the Cardano ledger, which is beyond the scope of
this specification (but see~\cite{plinth-userguide:contexts} for more information).


\begin{table}[H]
  \centering
    \begin{tabular}{|c|c|}
        \hline
        \thead{Protocol version} & \thead{Ledger language introduced} \\
        \hline
        5.0  & \LL{PlutusV1} \\
        7.0  & \LL{PlutusV2} \\
        9.0  & \LL{PlutusV3} \\
        \hline
        12.0 & \LL{PlutusV4} \\
        \hline
    \end{tabular}
    \caption{Introduction of Plutus ledger languages}
    \label{table:ll-introduction}
\end{table}

\subsection{Evolution of Plutus ledger languages and Plutus Core language versions}
Plutus Core scripts first became available on Cardano at the Alonzo hard fork
(PV5).  Later hard forks have introduced a number of ``batches'' of built-in
type and functions (described in Section~\ref{sec:builtin-batches} below) and
new Plutus Core language versions (see Section~\ref{sec:grammar-notes}).
Initially the workings of the Cardano ledger meant that it was not possible to
retroactively add new built-in functions or Plutus Core language versions to
existing ledger languages after a new ledger language was introduced: for
example Batch 3 was added to PlutusV2 at the Valentine hard fork (introducing
PV8), but it was not possible to add Batch 3 to PlutusV1 at that time.  This
restriction was subsequently lifted and at the van Rossem hard fork (PV11) it
became possible to use all built-in types and functions and all Plutus Core
language versions in all ledger languages.  Table~\ref{table:ll-pv-features}
below shows how the situation has evolved as time progressed (the row for PV12
is provisional at the time of writing).  This is of some relevance for nodes
replaying the history of the chain because they must treat scripts which use
builtins or Plutus Core versions which were not available at the time when the
scripts were originally executed as failing.\footnote{In fact, in PV7 and later,
scripts which use an unavailable ledger language or unavailable built-in types
or functions should be rejected during deserialisation and should never be
recorded on the chain.  These are \textit{Phase 1 failures}, as opposed to
\textit{Phase 2 failures}, which involve run-time errors such as division by
zero or budget overflows, and also attempts to use an unavailable Plutus Core
version.  Before PV7 scripts were only deserialised when they were executed, so
scripts using unavailable features would have caused Phase 2 failures instead.
Script evaluations with Phase 2 failures will be recorded in the history of the
chain, although examination of the history shows that very few Phase 2 failures
have actually occurred (the deterministic nature of script evaluation means that
they become apparent if a script is evaluated off-chain prior to submission, so
it is unlikely that scripts with Phase 2 failures will be submitted to the
chain).}


\begin{table}[H]
  \centering
  \small
  \setlength{\tabcolsep}{4pt}
  \renewcommand{\cellalign}{cc}
  \begin{tabular}{|l|c|c|c|c|}
    \hline
    \thead{Protocol\\version}
      & \thead{\LL{PlutusV1}} & \thead{\LL{PlutusV2}}
      & \thead{\LL{PlutusV3}} & \thead{\LL{PlutusV4}} \\
    \hline
    5.0 (Alonzo)
      & \makecell{1.0.0\\[1pt] Batch 1}
      & --- & --- & --- \\
    \hline
    7.0 (Vasil)
      & \makecell{1.0.0\\[1pt] Batch 1}
      & \makecell{Batches  1.0.0\\[1pt] 1--2}
      & --- & --- \\
    \hline
    8.0 (Valentine)
      & \makecell{1.0.0\\[1pt] Batch 1}
      & \makecell{1.0.0\\[1pt] Batches 1--3}
      & --- & --- \\
    \hline
    9.0 (Chang)
      & \makecell{1.0.0\\[1pt] Batch 1}
      & \makecell{1.0.0\\[1pt] Batches 1--3}
      & \makecell{1.0.0, 1.1.0 \\[1pt] Batches 1--4}
      & --- \\
    \hline
    10.0 (Plomin)
      & \makecell{1.0.0\\[1pt] Batch 1}
      & \makecell{1.0.0\\[1pt] Batches 1--3, 4b$^{\dagger}$}
      & \makecell{1.0.0, 1.1.0\\[1pt] Batches 1--5}
      & --- \\
    \hline
    11.0 (van Rossem)
      & \makecell{1.0.0, 1.1.0\\[1pt] Batches 1--6}
      & \makecell{1.0.0, 1.1.0\\[1pt] Batches 1--6}
      & \makecell{1.0.0, 1.1.0\\[1pt] Batches 1--6}
      & --- \\
    \hline
    \hline
    12.0 (Dijkstra HF)
      & \makecell{1.0.0--1.2.0\\[1pt] Batches 1--7}
      & \makecell{1.0.0--1.2.0\\[1pt] Batches 1--7}
      & \makecell{1.0.0--1.2.0\\[1pt] Batches 1--7}
      & \makecell{1.0.0--1.2.0\\[1pt] Batches 1--7}\\
    \hline
  \end{tabular}
  \caption{Built-in functions and Plutus Core language versions available in
    each ledger language at each major protocol version}
  \label{table:ll-pv-features}
\end{table}

\noindent Each entry in the table shows the Plutus Core language versions which
are allowed and the built-in batches which are available (see
Section~\ref{sec:builtin-batches}).  An entry of `---' means that the
ledger language does not exist at that protocol version.  The entry for PV12 is
provisional at the time of writing.  Only protocol versions at which something
changes are shown: for example, no relevant changes happened at PV6.  At PV10,
the Batch 4b built-ins (\TT{integerToByteString} and \TT{byteStringToInteger})
were added to \LL{PlutusV2} as a proof of concept; the remainder of Batch 4
became available in \LL{PlutusV2} at PV11.

\section{Built-in types and functions}
\label{sec:cardano-builtins}
\paragraph{Built-in batches.}
\label{sec:builtin-batches}

The built-in types and functions are defined in batches corresponding to how
they were added to ledger languages.  These batches are described in this section,
after some preliminaries.

\paragraph{Built-in semantics variants.}
\label{sec:builtin-semantics-variants}

In rare cases we can make a mistake or need to change the actual behaviour of a
built-in function.  To handle this we define a series of built-in semantics
variants, which indicate which behaviour should be used.  A fix will typically
be deployed by defining a new semantics variant, and then using that variant for
future ledger languages (but not existing ones, since this is usually a
backwards-incompatible change).  Table~\ref{table:bs-introduction} shows which
semantics variants are used for which ledger langauges.

\begin{table}[H]
  \centering
    \begin{tabular}{|c|c|c|}
        \hline
        \thead{Ledger language} & \thead{Built-in semantics variant used} \\
        \hline
        \LL{PlutusV1}, \LL{PlutusV2} & Built-in semantics 1 \\
        \LL{PlutusV3}, \LL{PlutusV4} & Built-in semantics 2 \\
        \hline
    \end{tabular}
    \caption{Selection of built-in semantics variant}
    \label{table:bs-introduction}
\end{table}

\noindent Table~\ref{table:bs-variants} shows the semantic changes for the
different variants.  Changes are listed alongside the original definition of the
built-in function in its original batch.

\begin{table}[H]
  \centering
    \begin{tabular}{|l|l|l|}
        \hline
        \thead{Built-in semantics variant} & \thead{Changes from previous semantics} \\
        \hline
        Built-in semantics 1 & None \\
        Built-in semantics 2 & \TT{consByteString} (see Note~\ref{note:consbytestring} in Section~\ref{sec:built-in-functions-1}) \\
        \hline
    \end{tabular}
    \caption{Built-in semantics variants}
    \label{table:bs-variants}
\end{table}


\paragraph{Concrete syntax for built-in types.}
Recall that in the abstract notation for built-in types introduced in
Section~\ref{sec:built-in-types}, a built-in type is either an \textit{atomic
  type} such as \texttt{integer} or \texttt{string} or an application $\op(T_1,
\ldots, T_n)$ of a \textit{type operator} to a sequence of built-in types.  The
concrete syntax of built-in types used in textual Plutus Core programs is
slightly different in that we use a curried form of application for type
operators: a type is given by

\newcommand{\bitn}{\mathbf{T}} % built-in type name
\newcommand{\bitc}{\mathbf{C}} % built-in constant syntax

\begin{minipage}{\linewidth}
    \centering
    \[\begin{array}{rclr}
    \bitn & ::= & \textit{atomic-type} & \textrm{Atomic type}\\
                & & \texttt{(}\op \ \bitn_1 \ldots \bitn_{\valency{\op}}) & \textrm{Type application}\\
    \end{array}\]
    \label{fig:built-in-type-concrete-syntax}
\end{minipage}
\noindent Note that we again require that all type operators are fully applied.
We refer to the syntactic objects $\bitn$ above as \textit{concrete built-in
 types}. There is an obvious bijection between these and the abstract built-in
types used elsewhere in this document, and given an abstract built-in type $T$
we will denote the corresponding concrete built-in type by $\bar{T}$.

\paragraph{Concrete syntax for built-in constants.}
We provide concrete syntax for constants of most (but not all) built-in types.
For a built-in type $T$ which has a concrete syntax we specify a set $\bitc_T$
of strings (using either regular expressions or a BNF-style grammar), and a
constant of type $T$ is then represented in the concrete syntax by an expression
of the form \texttt{(con $\bar{T}$ $c_T$)} with $c_T \in \bitc_T$.  Each string
$c_T$ will have an interpretation as a value of type $T$ (ie, an element of
$\denote{T}$) and since this will generally be the obvious interpretation we
will not always spell out the details.%%
\nomenclature[H]{$\bitc_T$}{Set of strings used for the concrete syntax of constants of built-in type $T$}

\newcounter{notenumberA}
\newcommand{\note}[1]{
  \bigskip
  \refstepcounter{notenumberA}
  \noindent\textbf{Note \thenotenumberA. #1}
}

\newcommand{\utfeight}{\mathsf{utf8}}
\newcommand{\unutfeight}{\mathsf{utf8}^{-1}}
\newcommand{\vk}{\textit{vk}}  %% Verification key (ie public key) for signature verification functions.

\subsection{Batch 1}
\label{sec:default-builtins-1}

\subsubsection{Built-in types and type operators}
\label{sec:built-in-types-1}
The first batch of built-in types and type operators is defined in Tables~\ref{table:built-in-types-1}
and~\ref{table:built-in-type-operators-1}.  We also include concrete syntax for
these; the concrete syntax is not strictly part of the language, but may be
useful for tools working with Plutus Core.

\begin{table}[H]
  \centering
    \begin{tabular}{|l|p{6cm}|l|}
        \hline
        Type & Denotation & Concrete Syntax\\
        \hline
        \texttt{integer} &   $\mathbb{Z}$ & \texttt{-?[0-9]+}\\
        \texttt{bytestring}  & $ \B^*$, the set of sequences of bytes or 8-bit characters. & \texttt{\#([0-9A-Fa-f][0-9A-Fa-f])*}\\
        \texttt{string} & $\U^*$,  the set of sequences of Unicode characters. & See note below\\
        \texttt{bool} & \{\true, \false\} & \texttt{True | False}\\
        \texttt{unit} &  \{()\} & \texttt{()}\\
        \texttt{data} &  See below & See below\\
        \hline
    \end{tabular}
    \caption{Atomic types, Batch 1}
    \label{table:built-in-types-1}
\end{table}

\begin{table}[H]
  \centering
    \begin{tabular}{|l|p{14mm}|l|l|}
        \hline
        Operator $\op$ & $\left|\op\right|$  & Denotation & Concrete Syntax\\
        \hline
        \texttt{list} & 1 & $\denote{\listOf{t}} = \denote{t}^*$ & See below\\
        \texttt{pair} & 2 & $\denote{\pairOf{t_1}{t_2}} = \denote{t_1} \times \denote{t_2}$ & See below\\
        \hline
        \end{tabular}
   \caption{Type operators, Batch 1}
    \label{table:built-in-type-operators-1}
\end{table}

\paragraph{Casing on built-in constants.}
From Protocol Version 11 onwards, Plutus Core supports casing on some of the
built-in types described in Tables~\ref{table:built-in-types-1} and
\ref{table:built-in-type-operators-1}.  See Table~\ref{table:builtin-matching-functions}
and Section~\ref{sec:builtin-casing} for details.  Casing is only available for
the types mentioned in Table~\ref{table:builtin-matching-functions}, and for all
other types $T$ one should assume that $\maxbranch(T) = 0$ and $\match{T}(x)
= \errorX$ for all $x$: the semantics of Plutus Core as defined in
Figure~\ref{fig:untyped-term-reduction} and implemented in the CEK machine
(Figure~\ref{fig:untyped-cek-machine}) will ensure that this
causes \texttt{(error)} if an attempt is made to case on a value of such a type
during evaluation of a Plutus Core script.

{ % New scope: change the vertical spacing for this table only.
% Casing functions for individual types are represented as embedded arrays.
% We still use an array even if there's only one case: without this the
% sub-arrays are indented further than the non-arrays.
\renewcommand*{\arraystretch}{2.3}
\begin{table}[H]
  \centering
    \begin{tabular}{|l|l|l|}
        \hline
        Type $T$ &\quad $\match{T}$  & $\maxbranch(T)$\\
        \hline
% unit
        $\ty{unit}$ &
        {\renewcommand*{\arraystretch}{1}   % Reset spacing for sub-array
        $ \begin{array}{ll}
         \mathtt{()} &\mapsto (1,())\\
        \end{array}$
        }
        & $1$
        \\
% boolean
        $\ty{boolean}$ &
        {\renewcommand*{\arraystretch}{1}
        $ \begin{array}{ll}
         \mathtt{false} & \mapsto (1,())\\
         \mathtt{true}  &\mapsto (2,())
        \end{array}$
        }
        & $2$
        \\
% integer
        \ty{integer}
        &
        {\renewcommand*{\arraystretch}{1}
        % We have k+1 here because we're numbering the branches from 1 to n, not 0 to n-1;
        % negative integers have to yield \errorX because \match{T} is required to return
        % an index in \Nplus.
        $\begin{array}{lll}
         k &\mapsto (k+1,()) & (k \geq 0)\\
         k &\mapsto \errorX  & (k<0)
        \end{array}$
        }
        & $\infty$
        \\
% list
        $\ty{list}\,T_1$ &
        {\renewcommand*{\arraystretch}{1}
        $ \begin{array}{ll}
         [x_1,x_2,\ldots,x_n] & \mapsto (1,(x_1,[x_2,\ldots,x_n]))   \quad (n \geq 1)\\
         %% \\ followed by [] causes problems because it looks like a size argument for \\
         \left[ \right]                  & \mapsto (2,())\\
        \end{array}$
        }
        & $2$
        \\
% pair
        $\ty{pair}\, T_1\, T_2$ &
        {\renewcommand*{\arraystretch}{1}
        $ \begin{array}{ll}
        (a,b) & \mapsto (1,(a,b))\\
        \end{array}$
        }
        & $1$
        \\
        \hline
        \end{tabular}
   \caption{Matching functions for built-in types}
    \label{table:builtin-matching-functions}
\end{table}
}



\paragraph{Concrete syntax for strings.} Strings are represented as sequences of Unicode characters
enclosed in double quotes, and may include standard escape sequences.  Surrogate
characters in the range \texttt{U+D800}--\texttt{U+DFFF} are replaced with the
Unicode replacement character \texttt{U+FFFD}.


\paragraph{Concrete syntax for lists and pairs.}
A list of type $\texttt{list}(t)$ is written as a syntactic list
\texttt{[$c_1, \ldots, c_n$]} where each $c_i$ lies in $\bitc_t$; a pair of type
$\texttt{pair}(t_1,t_2)$ is written as a syntactic pair $\texttt{(}c_1,c_2\texttt{)}$
with $c_1 \in \bitc_{t_1}$ and $c_2 \in \bitc_{t_2}$.  Some valid constant expressions
are thus

\begin{verbatim}
   (con (list integer) [11, 22, 33])
   (con (pair bool string) (True, "Plutus")).
   (con (list (pair bool (list bytestring)))
      [(True, []), (False, [#,#1F]), (True, [#123456, #AB, #ef2804])])
\end{verbatim}


\paragraph{The $\ty{data}$ type.}
We provide a built-in type $\ty{data}$ which permits the encoding of simple data
structures for use as arguments to Plutus Core scripts.  This type is defined in
Haskell as
\begin{alltt}
   data Data =
      Constr Integer [Data]
      | Map [(Data, Data)]
      | List [Data]
      | I Integer
      | B ByteString
\end{alltt}

\noindent In set-theoretic terms the denotation of $\ty{data}$ is
defined to be the least fixed point of the endofunctor $F$ on the category of
sets given by $F(X) = (\denote{\ty{integer}} \times X^*) \disj (X \times X)^* \disj
X^* \disj \denote{\ty{integer}} \disj \denote{\ty{bytestring}}$, so that
$$ \denote{\ty{data}} = (\denote{\ty{integer}} \times \denote{\ty{data}}^*)
               \disj (\denote{\ty{data}} \times \denote{\ty{data}})^*
               \disj \denote{\ty{data}}^*
               \disj \denote{\ty{integer}}
               \disj \denote{\ty{bytestring}}.
$$
We have injections
\begin{align*}
  \inj_C: \denote{\ty{integer}} \times \denote{\ty{data}}^* & \to \denote{\ty{data}} \\
  \inj_M: (\denote{\ty{data}} \times \denote{\ty{data}})^*  & \to \denote{\ty{data}} \\
  \inj_L: \denote{\ty{data}}^* & \to \denote{\ty{data}} \\
  \inj_I: \denote{\ty{integer}} & \to \denote{\ty{data}} \\
  inj_B: \denote{\ty{bytestring}} & \to \denote{\ty{data}} \\
\end{align*}
\noindent and projections
\begin{align*}
  \proj_C: \denote{\ty{data}} & \to \withError{(\denote{\ty{integer}} \times \denote{\ty{data}}^*)}\\
  \proj_M: \denote{\ty{data}} & \to \withError{(\denote{\ty{data}} \times \denote{\ty{data}})^*}\\
  \proj_L: \denote{\ty{data}} & \to \withError{\denote{\ty{data}}^* }\\
  \proj_I: \denote{\ty{data}} & \to \withError{\denote{\ty{integer}}}\\
  \proj_B: \denote{\ty{data}} & \to \withError{\denote{\ty{bytestring}} }\\
\end{align*}
\noindent which extract an object of the relevant type from a $\ty{data}$ object
$D$, returning $\errorX$ if $D$ does not lie in the expected component of the
disjoint union; also there are functions
$$
\is_C, \is_M, \is_L, \is_I, \is_B: \denote{\ty{data}} \to \denote{\ty{bool}}
$$
\noindent which determine whether a $\ty{data}$ value lies in the relevant component.

\paragraph{Note: \texttt{Constr} tag values.}
\label{note:constr-tag-values}
The \texttt{Constr} constructor of the \texttt{data} type is intended to
represent values from algebraic data types (also known as sum types and
discriminated unions, among other things; \texttt{data} itself is an example of
such a type), where $\mathtt{Constr}\, i\, [d_1,\ldots,d_n]$
represents a tuple of data items together with a tag $i$ indicating which of a
number of alternatives the data belongs to.  The definition above allows tags to
be any integer value, but because of restrictions in the serialisation format
for \texttt{data} (see Section~\ref{sec:encoding-data}) we recommend that in
practice \textbf{only tags $i$ with $0 \leq i \leq 2^{64}-1$ should be used}:
deserialisation will fail for \texttt{data} items (and programs which include
such items) involving tags outside this range.

Note also that \texttt{Constr} is unrelated to the $\keyword{constr}$ term in
Plutus Core itself. Both provide ways of representing structured data, but
the former is part of a built-in type whereas the latter is part of the language
itself.

\newcommand{\syn}[1]{c_{\mathtt{{#1}}}}

\paragraph{Concrete syntax for $\ty{data}$.}
The concrete syntax for $\ty{data}$ is given by

\begin{minipage}{0.6\linewidth}
    \centering
    \[\begin{array}{rcl}
    \syn{data} & ::= & \texttt{(Constr} \ \syn{integer} \ \syn{list(data)} \texttt{)}\\
               &  & \texttt{(Map} \ \syn{list(pair(data,data))} \texttt{)}\\
               &  & \texttt{(List} \ \syn{list(data)} \texttt{)}\\
               &  & \texttt{(I} \ \syn{integer} \texttt{)}\\
               &  & \texttt{(B} \ \syn{bytestring} \texttt{)}.
    \end{array}\]
    \label{fig:data-concrete-syntax}
\end{minipage}

\noindent We interpret these syntactic constants as elements of $\denote{\ty{data}}$ using
the various `$\inj$' functions defined earlier.  Some valid \texttt{data}
constants are

\begin{verbatim}
   (con data (Constr 1 [(I 2), (B #), (Map [])])
   (con data (Map [((I 0), (B #00)), ((I 1), (B #0F))]))
   (con data (List [(I 0), (I 1), (B #7FFF), (List []])))
   (con data (I -22))
   (con data (B #001A)).
\end{verbatim}
% TODO: be more relaxed about parenthesisation in general

\paragraph{Note.}  At the time of writing the syntax accepted by IOG's parser for textual Plutus Core
differs slightly from that above in that subobjects
of \texttt{Constr}, \texttt{Map} and \texttt{List} objects must \textit{not} be
parenthesised: for example one must write \verb|(con data (Constr 1 [I 2, B #,Map []])|.
This discrepancy will be resolved in the near future.


\subsubsection{Built-in functions}
\label{sec:built-in-functions-1}
The first batch of built-in functions is shown in
Table~\ref{table:built-in-functions-1}.  The table indicates which functions can
fail during execution, and conditions causing failure are specified either in
the denotation given in the table or in a relevant note.  Recall also that a
built-in function will fail if it is given an argument of the wrong type: this
is checked in conditions involving the $\sim$ relation and the $\EvalBuiltin$ function
in Figures~\ref{fig:untyped-term-reduction} and~\ref{fig:untyped-cek-machine}.
Note also that some of the functions are
\#-polymorphic.  According to Section~\ref{sec:builtin-denotations} we
require a denotation for every possible monomorphisation of these; however all
of these functions are parametrically polymorphic so to simplify notation we
have given a single denotation for each of them with an implicit assumption that
it applies at each possible monomorphisation in an obvious way.

\setlength{\LTleft}{-18mm} % Shift the table left a bit to centre it on the page
\begin{longtable}[H]{|l|p{5cm}|p{5.5cm}|c|c|}
    \hline
    \text{Function} & \text{Signature} & \text{Denotation} & \text{Can} & \text{Note} \\
    & & & fail? & \\
    \hline
    \endfirsthead
    \hline
    \text{Function} & \text{Type} & \text{Denotation} & \text{Can} & \text{Note}\\
    & & & fail? & \\
    \hline
    \endhead
    \hline
    \caption{Built-in functions, Batch 1}
    % This caption goes on every page of the table except the last.  Ideally it
    % would appear only on the first page and all the rest would say
    % (continued). Unfortunately it doesn't seem to be easy to do that in a
    % longtable.
    \endfoot
    \caption[]{Built-in functions, Batch 1 (continued)}
    \label{table:built-in-functions-1}
    \endlastfoot
    \TT{addInteger}               & $[\ty{integer}, \ty{integer}] \to \ty{integer}$   & $+$ & No & \\[2mm]
    \TT{subtractInteger}          & $[\ty{integer}, \ty{integer}] \to \ty{integer}$   & $-$ & No & \\[2mm]
    \TT{multiplyInteger}          & $[\ty{integer}, \ty{integer}] \to \ty{integer}$   & $\times$ & No & \\[2mm]
    \TT{divideInteger}            & $[\ty{integer}, \ty{integer}] \to \ty{integer}$   & $\divfn$   & Yes & \ref{note:integer-division-functions}\\[2mm]
    \TT{modInteger}               & $[\ty{integer}, \ty{integer}] \to \ty{integer}$   & $\modfn$   & Yes & \ref{note:integer-division-functions}\\[2mm]
    \TT{quotientInteger}          & $[\ty{integer}, \ty{integer}] \to \ty{integer}$   & $\quotfn$  & Yes & \ref{note:integer-division-functions}\\[2mm]
    \TT{remainderInteger}         & $[\ty{integer}, \ty{integer}] \to \ty{integer}$   & $\remfn$   & Yes & \ref{note:integer-division-functions}\\[2mm]
    \TT{equalsInteger}            & $[\ty{integer}, \ty{integer}] \to \ty{bool}$      & $=$ & No & \\[2mm]
    \TT{lessThanInteger}          & $[\ty{integer}, \ty{integer}] \to \ty{bool}$      & $<$ & No & \\[2mm]
    \TT{lessThanEqualsInteger}    & $[\ty{integer}, \ty{integer}] \to \ty{bool}$      & $\leq$ & No & \\[2mm]
    %% Some of the signatures look like $ ... $ \text{\;\; $ ... $} to allow a break with some indentation afterwards
    \TT{appendByteString}         & $[\ty{bytestring}, \ty{bytestring}] $ \text{$\;\; \to \ty{bytestring}$}
                                           & $([c_1, \dots, c_m], [d_1, \ldots, d_n]) $ \text{$\;\; \mapsto [c_1,\ldots, c_m,d_1, \ldots, d_n]$} & No & \\[2mm]
    \TT{consByteString} (Variant 1) & $[\ty{integer}, \ty{bytestring}] $ \text{$\;\; \to \ty{bytestring}$}
                                          & $(c,[c_1,\ldots,c_n]) $ \text{$\;\;\mapsto [\text{mod}(c,256) ,c_1,\ldots,c_{n}]$} & No
                                          & \ref{note:consbytestring}\\[2mm]
    \TT{consByteString} (Variant 2) & $[\ty{integer}, \ty{bytestring}] $ \text{$\;\; \to \ty{bytestring}$}
                                          & $(c,[c_1,\ldots,c_n])$ \text{$\;\;\mapsto
                                                       \begin{cases}
                                                          [c,c_1,\ldots,c_{n}] & \text{if $0 \leq c \leq 255$} \\[2mm]
                                                         \errorX & \text{otherwise}
                                                       \end{cases}$} & Yes & \ref{note:consbytestring}\\[2mm]
    \TT{sliceByteString}        & $[\ty{integer}, \ty{integer}, \ty{bytestring]} $  \text {$\;\; \to  \ty{bytestring}$}
                                                   &   $(s,k,[c_1,\ldots,c_n])$ \text{$\;\;\mapsto [c_{\max(s+1,1)},\ldots,c_{\min(s+k,n)}]$}
                                                   & No & \ref{note:slicebytestring}\\[2mm]
    \TT{lengthOfByteString}       & $[\ty{bytestring}] \to \ty{integer}$
            & $[] \mapsto 0, [c_1,\ldots, c_n] \mapsto n \ (n \geq 1)$ & No & \\[2mm]
    \TT{indexByteString}          & $[\ty{bytestring}, \ty{integer}] $ \text{$\;\; \to \ty{integer}$}
                                                   & $([c_1,\ldots,c_n],i)$ \text{$\;\;\mapsto
                                                       \begin{cases}
                                                         c_{i+1} & \text{if $0 \leq i \leq n-1$} \\[2mm]
                                                         \errorX & \text{otherwise}
                                                       \end{cases}$} & Yes & \\[2mm]
    \TT{equalsByteString}         & $[\ty{bytestring}, \ty{bytestring}] $ \text{$\;\; \to \ty{bool}$}   & = & No & \ref{note:bytestring-comparison}\\[2mm]
    \TT{lessThanByteString}       & $[\ty{bytestring}, \ty{bytestring}] $ \text{$\;\; \to \ty{bool}$}   & $<$ & No & \ref{note:bytestring-comparison}\\[2mm]
    \TT{lessThanEqualsByteString} & $[\ty{bytestring}, \ty{bytestring}] $ \text{$\;\; \to \ty{bool}$}   & $\leq$ & No & \ref{note:bytestring-comparison}\\[2mm]
    \TT{appendString}             & $[\ty{string}, \ty{string}] \to \ty{string}$
                                         & $([u_1, \dots, u_m], [v_1, \ldots, v_n]) $ \text{$\;\; \mapsto [u_1,\ldots, u_m,v_1, \ldots, v_n]$} & No & \\[2mm]
    \TT{equalsString}             & $[\ty{string}, \ty{string}] \to \ty{bool}$           & = & No & \\[2mm]
    \TT{encodeUtf8}               & $[\ty{string}] \to \ty{bytestring}$      & $\utfeight$ & & \ref{note:bytestring-encoding} \\[2mm]
    \TT{decodeUtf8}               & $[\ty{bytestring}] \to \ty{string}$      & $\unutfeight$ & Yes & \ref{note:bytestring-encoding} \\[2mm]
    \TT{sha2\_256}                & $[\ty{bytestring}] \to \ty{bytestring}$  & \text{Hash a $\ty{bytestring}$ using} \TT{SHA-}\TT{256}~\cite{FIPS-SHA2}. & No & \\[2mm]
    \TT{sha3\_256}                & $[\ty{bytestring}] \to \ty{bytestring}$  & \text{Hash a $\ty{bytestring}$ using} \TT{SHA3-}\TT{256}~\cite{FIPS-SHA3}. & No & \\[2mm]
    \TT{blake2b\_256}             & $[\ty{bytestring}] \to \ty{bytestring}$  & \text{Hash a $\ty{bytestring}$ using} \TT{Blake2b-}\TT{256}~\cite{IETF-Blake2}. & No & \\[2mm]
    \TT{verifyEd25519Signature}          & $[\ty{bytestring}, \ty{bytestring}, $ \text{$\;\; \ty{bytestring}] \to \ty{bool}$}
    & Verify an \TT{Ed25519} digital signature. &  Yes
    & \ref{note:digital-signature-verification-functions}, \ref{note:ed25519-signature-verification}\\[2mm]
    \TT{ifThenElse}               & $[\forall\star, \ty{bool}, \star, \star] \to \star$
                                                 & \text{$(\true,t_1,t_2) \mapsto t_1$}
                                                 \text{$(\false,t_1,t_2) \mapsto t_2$} & No & \\[2mm]
    \TT{chooseUnit}               & $[\forall\star, \ty{unit}, \star] \to \star$        & $((), t) \mapsto t$ & No & \\[2mm]
    \TT{trace}                    & $[\forall\star, \ty{string}, \star] \to \star$      & $ (s,t) \mapsto t$ & No & \ref{note:trace}\\[2mm]
    \TT{fstPair}                  & $[\forall a_\#, \forall b_\#, \pairOf{a_\#}{b_\#}] \to a_\#$       & $(x,y) \mapsto x$ & No & \\[2mm]
    \TT{sndPair}                  & $[\forall a_\#, \forall b_\#, \pairOf{a_\#}{b_\#}] \to b_\#$       & $(x,y) \mapsto y$ & No & \\[2mm]
    \TT{chooseList}               & $[\forall a_\#, \forall\star, \listOf{a_\#}, \star, \star] \to \star$
                                              & \text{$([], t_1, t_2) \mapsto t_1$,} \text{$([x_1,\ldots,x_n],t_1,t_2) \mapsto t_2\ (n \geq 1)$}. & No & \ref{note:list-semantics}\\[2mm]
    \TT{mkCons}                   & $[\forall a_\#, a_\#, \listOf{a_\#}] \to \listOf{a _\#}$  & $(x,[x_1,\ldots,x_n]) \mapsto [x,x_1,\ldots,x_n]$ & No & \ref{note:list-semantics}\\[2mm]
    \TT{headList}                 & $[\forall a_\#, \listOf{a_\#}] \to a_\#$               & $[]\mapsto \errorX, [x_1,x_2, \ldots, x_n] \mapsto x_1$ & Yes & \ref{note:list-semantics}\\[2mm]
    \TT{tailList}                 & $[\forall a_\#, \listOf{a_\#}] \to \listOf{a_\#}$
                                        &  \text{$[] \mapsto \errorX$,} \text{$ [x_1,x_2, \ldots, x_n] \mapsto [x_2, \ldots, x_n]$} & Yes & \ref{note:list-semantics}\\[2mm]
    \TT{nullList}                 & $[\forall a_\#, \listOf{a_\#}] \to \ty{bool}$            & $ [] \mapsto \true,
                                                                                                    [x_1,\ldots, x_n] \mapsto \false$ & No & \ref{note:list-semantics} \\[2mm]
    \TT{chooseData}               & $[\forall\star, \ty{data}, \star, \star, \star, \star, \star] \to \star$
    & $ (d,t_C, t_M, t_L, t_I, t_B) $
    \smallskip
    \newline  % The big \{ was abutting the text above
    \text{$\;\;\mapsto
               \left\{ \begin{array}{ll}  %% This looks better than `cases`
                 t_C  & \text{if $\is_C(d)$} \\
                 t_M  & \text{if $\is_M(d)$} \\
                 t_L  & \text{if $\is_L(d)$} \\
                 t_I  & \text{if $\is_I(d)$} \\
                 t_B  & \text{if $\is_B(d)$} \\
               \end{array}\right.$}  & No & \\
    \TT{constrData}               & $[\ty{integer}, \listOf{\ty{data}}] \to \ty{data}$          & $\inj_C$ & No & \\[2mm]
    \TT{mapData}                  & $[\listOf{\pairOf{\ty{data}}{\ty{data}}}$ \text{$\;\; \to \ty{data}$}     & $\inj_M$ & No & \\[2mm]
    \TT{listData}                 & $[\listOf{\ty{data}}] \to \ty{data} $      & $\inj_L$ & No & \\[2mm]
    \TT{iData}                    & $[\ty{integer}] \to \ty{data} $            & $\inj_I$ & No & \\[2mm]
    \TT{bData}                    & $[\ty{bytestring}] \to \ty{data} $         & $\inj_B$& No & \\[2mm]
    \TT{unConstrData}             & $[\ty{data}]$ \text{$\;\; \to \pairOf{\ty{integer}}{\listOf{\ty{data}}}$} & $\proj_C$ & Yes & \\[2mm]
    \TT{unMapData}                & $[\ty{data}]$ \text{$\;\; \to \listOf{\pairOf{\ty{data}}{\ty{data}}}$}  & $\proj_M$ & Yes & \\[2mm]
    \TT{unListData}               & $[\ty{data}] \to \listOf{\ty{data}} $                          & $\proj_L$ & Yes & \\[2mm]
    \TT{unIData}                  & $[\ty{data}] \to \ty{integer} $                                & $\proj_I$ & Yes & \\[2mm]
    \TT{unBData}                  & $[\ty{data}] \to \ty{bytestring} $                             & $\proj_B$ & Yes & \\[2mm]
    \TT{equalsData}               & $[\ty{data}, \ty{data}] \to \ty{bool} $                        & $ = $ & No & \\[2mm]
    \TT{mkPairData}               & $[\ty{data}, \ty{data}]$ \text{\;\; $\to \pairOf{\ty{data}}{\ty{data}}$}  & $(x,y) \mapsto (x,y) $ & No & \\[2mm]
    \TT{mkNilData}                & $[\ty{unit}] \to \listOf{\ty{data}} $                       & $() \mapsto []$ & No & \\[2mm]
    \TT{mkNilPairData}            & $[\ty{unit}] $ \text{$\;\; \to \listOf{\pairOf{\ty{data}}{\ty{data}}} $}   & $() \mapsto []$ & No & \\[2mm]
    \hline
\end{longtable}

\kwxm{Maybe try \texttt{tabulararray} to see what sort of output that gives for the big table.}

\note{Integer division functions.}
\label{note:integer-division-functions}
We provide four integer division functions: \texttt{divideInteger},
\texttt{modInteger}, \texttt{quotientInteger}, and \texttt{remainderInteger},
whose denotations are mathematical functions $\divfn, \modfn, \quotfn$, and
$\remfn$ which are modelled on the corresponding Haskell operations. Each of
these takes two arguments and will fail (returning $\errorX$) if the second one
is zero.  For all $a,b \in \Z$ with $b \ne 0$ we have
$$
\divfn(a,b) \times b + \modfn(a,b) = a
$$
$$
  |\modfn(a,b)| < |b|
$$\noindent and
$$
  \quotfn(a,b) \times b + \remfn(a,b) = a
$$
$$
  |\remfn(a,b)| < |b|.
$$
\nomenclature[Azz]{$\divfn$, $\modfn$}{Integer division operations}
\nomenclature[Azz]{$\quotfn$, $\remfn$}{Integer division operations}

\noindent The $\divfn$ and $\modfn$ functions form a pair, as do $\quotfn$ and $\remfn$;
$\divfn$ should not be used in combination with $\remfn$, nor should $\quotfn$ be used
with $\modfn$.

For positive divisors $b$, $\divfn$ truncates downwards and $\modfn$ always
returns a non-negative result ($0 \leq \modfn(a,b) \leq b-1$).  The $\quotfn$
function truncates towards zero.  Table~\ref{table:integer-division-signs} shows
how the signs of the outputs of the division functions depend on the signs of
the inputs; $+$ means $\geq 0$ and $-$ means $\leq 0$, but recall that for $b=0$
all of these functions return the error value $\errorX$.
\begin{table}[H]
  \centering
    \begin{tabular}{|cc|cc|cc|}
        \hline
        a & b & $\divfn$ & $\modfn$ & $\quotfn$ & $\remfn$ \\
        \hline
        $+$ & $+$ & $+$ & $+$ & $+$ & $+$ \\
        $-$ & $+$ & $-$ & $+$ & $-$ & $-$ \\
        $+$ & $-$ & $-$ & $-$ & $-$ & $+$ \\
        $-$ & $-$ & $+$ & $-$ & $+$ & $-$ \\
        \hline
        \end{tabular}
   \caption{Behaviour of integer division functions}
   \label{table:integer-division-signs}
\end{table}
%% -------------------------------
%% |   n  d | div mod | quot rem |
%% |-----------------------------|
%% |  41  5 |  8   1  |   8   1  |
%% | -41  5 | -9   4  |  -8  -1  |
%% |  41 -5 | -9  -4  |  -8   1  |
%% | -41 -5 |  8  -1  |   8  -1  |
%% -------------------------------

\note{The \texttt{consByteString} function.}
\label{note:consbytestring}
In built-in semantics variant 1, the first argument of \texttt{consByteString}
is an arbitrary integer which will be reduced modulo 256 before being prepended
to the second argument.  In built-in semantics variant 2 we require that the first
argument lies between 0 and 255 (inclusive): in any other case an error will
occur.

\note{The \texttt{sliceByteString} function.}
\label{note:slicebytestring}
The application \texttt{[[(builtin sliceByteString) (con integer $s$)] (con
    integer $k$)] (con bytestring $b$)]} returns the substring of $b$ of length
$k$ starting at position $s$; indexing is zero-based, so a call with $s=0$
returns a substring starting with the first element of $b$, $s=1$ returns a
substring starting with the second, and so on.  This function always succeeds,
even if the arguments are out of range: if $b=[c_1, \ldots, c_n]$ then the
  application above returns the substring $[c_i, \ldots, c_j]$ where
  $i=\max(s+1,1)$ and $j=\min(s+k, n)$; if $j<i$ then the empty string is returned.

\note{Comparisons of bytestrings.}
\label{note:bytestring-comparison}
Bytestrings are ordered lexicographically in the usual way. If we have $a =
  [a_1, \ldots, a_m]$ and $b = [b_1, \ldots, b_n]$ then (recalling that if $m=0$
  then $a=[]$, and similarly for $b$),
\begin{itemize}
\item $a = b$ if and only if $m=n$ and $a_i = b_i$ for $1 \leq i \leq m$.

\item $a \leq b$ if and only if one of the following holds:
\begin{itemize}
  \item $a = []$
  \item $m,n > 0$ and $a_1 < b_1$
  \item $m,n > 0$ and $a_1 = b_1$ and $[a_2,\ldots,a_m] \leq [b_2,\ldots,b_n]$.
\end{itemize}
\item $a < b$ if and only if $a \leq b$ and $a \neq b$.
\end{itemize}
\noindent For example, $\mathtt{\#23456789} < \mathtt{\#24}$ and
$\mathtt{\#2345} < \mathtt{\#234500}$.  The empty bytestring is equal only to
itself and is strictly less than all other bytestrings.

\note{Encoding and decoding bytestrings.}
\label{note:bytestring-encoding}
The \texttt{encodeUtf8} and \texttt{decodeUtf8} functions convert between the
$\ty{string}$ type and the $\ty{bytestring}$ type.  We have defined
$\denote{\ty{string}}$ to consist of sequences of Unicode characters without
specifying any particular character representation, whereas
$\denote{\ty{bytestring}}$ consists of sequences of 8-bit bytes.  We define the
denotation of \texttt{encodeUtf8} to be the function
$$
\utfeight: \U^* \rightarrow \B^*
$$

\noindent which converts sequences of Unicode characters to sequences of bytes using the
well-known UTF-8 character encoding~\cite[Definition  D92]{Unicode-standard}.
The denotation of \texttt{decodeUtf8} is the partial inverse function

$$
\unutfeight: \B^* \rightarrow \U^*_{\errorX}.
$$

\noindent UTF-8 encodes Unicode characters encoded using between one and four
bytes: thus in general neither function will preserve the length of an object.
Moreover, not all sequences of bytes are valid representations of Unicode
characters, and \texttt{decodeUtf8} will fail if it receives an invalid input
(but \texttt{encodeUtf8} will always succeed).


\kwxm{In fact, strings are represented as sequences of UTF-16 characters, which
  use two or four bytes per character.  Do we need to mention that?  If we do,
  we'll need to be a little careful: there are sequences of 16-bit words that
  don't represent valid Unicode characters (for example, if the sequence uses
  surrogate codepoints improperly.  I don't think you can create a Haskell
  \texttt{Text} object (which is what our strings really are) that's invalid
  though.}


\note{Digital signature verification functions.}
\label{note:digital-signature-verification-functions}
We use a uniform interface for digital signature verification algorithms. A
digital signature verification function takes three bytestring arguments (in the
given order):
\begin{itemize}
  \item a public key $\vk$ (in this context $\vk$ is also known as a \textit{verification key})
  \item a message $m$
  \item a signature  $s$.
\end{itemize}
\noindent A signature verification function may require one
or more arguments to be well-formed in some sense (in particular an argument
may need to be of a specified length), and in this case the function will fail
(returning $\errorX$) if any argument is malformed. If all of the arguments are
well-formed then the verification function returns $\true$ if the private
key corresponding to $\vk$ was used to sign the message $m$ to produce $s$,
otherwise it returns $\false$.

\note{Ed25519 signature verification.}
\label{note:ed25519-signature-verification}
The \texttt{verifyEd25519Signature}
function performs cryptographic
signature verification using the Ed25519 scheme \cite{ches-2011-24091,
  rfc8032-EdDSA}, and conforms to the interface described in
Note~\ref{note:digital-signature-verification-functions}.  The arguments must
have the following sizes:
\begin{itemize}
\item $\vk$: 32 bytes
\item $m$: unrestricted
\item $s$: 64 bytes.
\end{itemize}

\note{The \texttt{trace} function.}
\label{note:trace}
An application \texttt{[(builtin trace) $s$ $v$]} ($s$ a \texttt{string}, $v$
any Plutus Core value) returns $v$.  We do not specify the semantics any
further.  An implementation may choose to discard $s$ or to perform some
side-effect such as writing it to a terminal or log file.

\note{Semantics of lists.}
\label{note:list-semantics}
The mathematical denotation of lists is as finite sequences of elements.  We
expect that in practice they will be implemented in such a way (for example as
linked lists)
that \texttt{chooseList}, \texttt{mkCons}, \texttt{headList}, \texttt{tailList},
and \texttt{nullList} are all constant-time operations.
\newpage

% I tried resetting the note number from V1-builtins here, but that made
% some hyperlinks wrong.  To get note numbers starting at one in each section, I
% think we have to define a new counter every time.
\newcounter{notenumberB}
\renewcommand{\note}[1]{
  \bigskip
  \refstepcounter{notenumberB}
  \noindent\textbf{Note \thenotenumberB. #1}
}

\subsection{Tranche 2}
\label{sec:default-builtins-2}

\subsubsection{Built-in functions}
\label{sec:built-in-functions-2}
The second tranche of builtin operations is defined in Table~\ref{table:built-in-functions-2}.

\setlength{\LTleft}{-10mm}  % Shift the table left a bit to centre it on the page
\begin{longtable}[H]{|l|p{42mm}|p{35mm}|c|c|}
    \hline
    \text{Function} & \text{Signature} & \text{Denotation} & \text{Can} & \text{Note} \\
    & & & fail? & \\
    \hline
    \endfirsthead
    \hline
    \text{Function} & \text{Type} & \text{Denotation} & \text{Can} & \text{Note}\\
    & & & fail? & \\
    \hline
    \endhead
    \hline
    \caption{Built-in Functions}
    \endfoot
    \caption[]{Built-in Functions}
    \label{table:built-in-functions-2}
    \endlastfoot
    \TT{serialiseData}                        & $[\ty{data}] \to \ty{bytestring}$   &  $\mathcal{E}_{\mathtt{data}}$ &
      & \ref{note:serialise-data}\\
\hline
\end{longtable}

\note{Serialising $\ty{data}$ objects.}
\label{note:serialise-data}
The \texttt{serialiseData} function takes a $\ty{data}$ object and converts it
into a bytestring using a CBOR encoding.  A full specification of the encoding
(including the definition of $\mathcal{E}_{\mathtt{data}}$) is provided in
Appendix~\ref{appendix:data-cbor-encoding}.

% I tried resetting the note number from V1-builtins here, but that made
% some hyperlinks wrong.  To get note numbers starting at one in each section, I
% think we have to define a new counter every time.
\newcounter{notenumberC}
\renewcommand{\note}[1]{
  \bigskip
  \refstepcounter{notenumberC}
  \noindent\textbf{Note \thenotenumberC. #1}
}

\subsection{Batch 3}
\label{sec:default-builtins-3}

\subsubsection{Built-in functions}
\label{sec:built-in-functions-3}
The third batch of built-in functions is defined in Table~\ref{table:built-in-functions-3}.
See~\cite{CIP-0049}.

\setlength{\LTleft}{-10mm}  % Shift the table left a bit to centre it on the page
\begin{longtable}[H]{|l|p{42mm}|p{35mm}|c|c|}
    \hline
    \text{Function} & \text{Signature} & \text{Denotation} & \text{Can} & \text{Note} \\
    & & & fail? & \\
    \hline
    \endfirsthead
    \hline
    \text{Function} & \text{Type} & \text{Denotation} & \text{Can} & \text{Note}\\
    & & & fail? & \\
    \hline
    \endhead
    \hline
    \caption{Built-in functions, Batch 3}
    \endfoot
    \caption[]{Built-in functions, Batch 3}
    \label{table:built-in-functions-3}
    \endlastfoot
    \TT{verifyEcdsaSecp256k1Signature}        & $[\ty{bytestring}, \ty{bytestring}, $ \text{$\;\; \ty{bytestring}] \to \ty{bool}$}
        & Verify an SECP-256k1 ECDSA signature & Yes & \ref{note:verify-ecdsa-secp256k1-signature}\\[2mm]
    \TT{verifySchnorrSecp256k1Signature}      & $[\ty{bytestring}, \ty{bytestring}, $ \text{$\;\; \ty{bytestring}] \to \ty{bool}$}
          & Verify an SECP-256k1 Schnorr signature & Yes & \ref{note:verify-schnorr-secp256k1-signature}\\[2mm]
\hline
\end{longtable}


\note{Secp256k1 ECDSA Signature verification.}
\label{note:verify-ecdsa-secp256k1-signature}
The \texttt{verifyEcdsaSecp256k1Signature} function performs elliptic curve
digital signature verification \cite{ANSI-X9.62, ANSI-x9.142,
  Johnson-Menezes-Vanstone-ECDSA} over the \texttt{secp256k1}
curve~\cite[\S2.4.1]{SECP256} and conforms to the interface described in
Note~\ref{note:digital-signature-verification-functions} of
Section~\ref{sec:built-in-functions-1}.  The arguments must have the
following sizes:
\begin{itemize}
\item $\vk$: 33 bytes
\item $m$: 32 bytes
\item $s$: 64 bytes.
\end{itemize}
The public key $\vk$ is expected to be in the 33-byte compressed form described
in~\cite{Bitcoin-ECDSA}.  It is expected that the inputs to the function are
well formed, and if they are not then an error will occur.  For example, the
signature $s$ encodes a pair $(R,S)$ of 32-byte bytestrings which are big-endian
encodings of integers, and both of these integers must lie between 0 and
$\left|G\right|-1$ (inclusive), where $G$ is the relevant subgroup of the
\texttt{secp256k1} curve; it is also required that the last 32 bytes of the
verification key are the big-endian encoding of the $x$-coordinate of a point in
$G$.

Moreover, the ECDSA scheme admits two distinct valid
signatures for a given message and private key, and  we follow the restriction
imposed by Bitcoin (see~\cite{BIP-146},
\texttt{LOW\_S}) and \textbf{only accept the smaller signature};
\texttt{verifyEcdsa\-Secp\-256k1Signature} will return $\false$ if the larger
one is supplied.

% For more on the lower signature business, see
%    https://github.com/IntersectMBO/plutus/pull/4591#issuecomment-1120797931
% and
%    https://github.com/IntersectMBO/plutus/pull/4591#issuecomment-1121684776
% The C code that performs the check in the bitcoin implementation is at
%   https://github.com/bitcoin-core/secp256k1/blob/485f608fa9e28f132f127df97136617645effe81/src/secp256k1.c#L400
% and
%   https://github.com/bitcoin-core/secp256k1/blob/485f608fa9e28f132f127df97136617645effe81/src/scalar_low_impl.h#L85
%
% Schnorr signatures do something substantially different for ECDSA and I don't think
% the question of multiple signatures arises for verifySchnorrSecp256k1Signature.

\note{Secp256k1 Schnorr Signature verification.}
\label{note:verify-schnorr-secp256k1-signature}
The \texttt{verifySchnorrSecp256k1Signature} function performs verification of
Schnorr signatures~\cite{Schnorr89, BIP-340} over the \texttt{secp256k1} curve
and conforms to the interface described in
Note~\ref{note:digital-signature-verification-functions} of
Section~\ref{sec:built-in-functions-1}.  The arguments are expected to be
of the forms specified in BIP-340~\cite{BIP-340} and thus should have the
following sizes:
\begin{itemize}
\item $\vk$: 32 bytes
\item $m$: unrestricted
\item $s$: 64 bytes.
\end{itemize}
\noindent  Any deviation from the conditions of BIP-340~\cite{BIP-340} will cause an error.

% I tried resetting the note number from V1-builtins here, but that made
% some hyperlinks wrong.  To get note numbers starting at one in each section, I
% think we have to define a new counter every time.
\newcounter{notenumberD}
\renewcommand{\note}[1]{
  \bigskip
  \refstepcounter{notenumberD}
  \noindent\textbf{Note \thenotenumberD. #1}
}

\newcommand{\itobsBE}{\mathsf{itobs_{BE}}}
\newcommand{\itobsLE}{\mathsf{itobs_{LE}}}
\newcommand{\bstoiBE}{\mathsf{bstoi_{BE}}}
\newcommand{\bstoiLE}{\mathsf{bstoi_{LE}}}

\subsection{Batch 4}
\label{sec:default-builtins-4}
The fourth batch of builtins adds support for
\begin{itemize}
\item The \texttt{Blake2b-224} and \texttt{Keccak-256} hash functions.
\item Conversions functions from integers to bytestrings and vice-versa.
\item BLS12-381 elliptic curve pairing operations
(see~\cite{CIP-0381}, \cite{BLS12-381}, \cite[4.2.1]{IETF-pairing-friendly-curves}, \cite{BLST-library}).
 For clarity these are described separately in Sections~\ref{sec:bls-types-4} and \ref{sec:bls-builtins-4}.
\end{itemize}

\subsubsection{Miscellaneous built-in functions}
\label{sec:misc-builtins-4}

\setlength{\LTleft}{-10mm} % Shift the table left a bit to centre it on the page
\begin{longtable}[H]{|l|p{45mm}|p{60mm}|c|c|}
    \hline \text{Function} & \text{Signature} & \text{Denotation} & \text{Can}
    & \text{Note} \\ & & & fail?
    & \\ \hline \endfirsthead \hline \text{Function} & \text{Type}
    % This caption goes on every page of the table except the last.  Ideally it
    % would appear only on the first page and all the rest would say
    % (continued). Unfortunately it doesn't seem to be easy to do that in a
    % longtable.
    \endfoot
%%    \caption[]{Built-in Functions}
    \caption[]{Miscellaneous built-in Functions}
    \label{table:misc-built-in-functions-4}
    \endlastfoot
%% G1
    \hline
    \TT{blake2b\_224} & $[\ty{bytestring}] \to \ty{bytestring}$  & \text{Hash a $\ty{bytestring}$ using}
                                                                   \TT{Blake2b-224}~\cite{IETF-Blake2}. &  & \\
    \TT{keccak\_256}  & $[\ty{bytestring}] \to \ty{bytestring}$  & \text{Hash a $\ty{bytestring}$ using}
                                                                   \TT{Keccak-256}~\cite{KeccakRef3}. &  & \\
    \hline\strut
    \TT{integerToByteString} & $[\ty{bool}, \ty{integer}, \ty{integer}]$  \text{\: $\to \ty{bytestring}$}
                                        & $(e, w, n) $ \text{$\mapsto \begin{cases}
                                        \itobsLE(w,n) & \text{if $e=\mathtt{false}$}\\
                                        \itobsBE(w,n) & \text{if $e=\mathtt{true}$}\\
                                        \end{cases}$}
                                        & Yes & \ref{note:itobs}\strut\\ \strut
    \TT{byteStringToInteger} & $[\ty{bool}, \ty{bytestring}] $ \text{\: $ \to \ty{bytestring}$}
                                        & $(e, [c_0, \ldots, c_{N-1}]) $ \text{\; $\mapsto \begin{cases}
                                        \sum_{i=0}^{N-1}c_i256^i & \text{if $e=\mathtt{false}$}\\
                                        \sum_{i=0}^{N-1}c_i256^{N-1-i} & \text{if $e=\mathtt{true}$}\\
                                        \end{cases}$}
                                        &  & \ref{note:bstoi}\\
    \hline
\end{longtable}

\note{Integer to bytestring conversion.}
\label{note:itobs}
The \texttt{integerToByteString} function converts non-negative integers to bytestrings.
It takes three arguments:
\begin{itemize}
\item An endianness flag $e$.
\item A width argument $w$.
\item The integer $n$ to be converted.
\end{itemize}

\noindent If $e$ is \texttt{True} then the conversion is big-endian, otherwise it is
little-endian. If the width $w$ is zero then the output is a bytestring which is
just large enough to hold the converted integer.  If $w>0$ then the output is
exactly $w$ bytes long, and it is an error if $n$ does not fit into a bytestring
of that size; if necessary, the output is padded with \texttt{0x00} bytes (on
the left if $e = \mathtt{True}$ and on the right if $e = \mathtt{False}$) to
make it the correct length.  In all cases it is an error if $n<0$.

\newpage
\noindent The precise semantics of \texttt{integerToByteString} are given
by the functions $\itobsBE: \Z \times \Z \rightarrow \B^*$ and $\itobsLE
: \Z \times \Z \rightarrow \B^*$.  Writing $n = \sum_{i=0}^{N-1}a_i256^i$ with
$a_{N-1} \ne 0$, we have

$$
\itobsLE (w,n) =
\begin{cases}
  [a_0, \ldots, a_{N-1}] & \text{if $w=0$} \\
  [b_0, \ldots, b_{w-1}] &  \text{if $w > 0$ and $N\leq w$, where }
      b_i = \begin{cases}
                a_i & \text{if $i \leq N-1$} \\
                0   & \text{if $i \geq N$}\\
            \end{cases}\\
  \errorX & \text{if $w > 0$ and $N > w$}
\end{cases}
$$

\noindent and

$$
\noindent
\itobsBE (w,n) =
\begin{cases}
  [b_0, \ldots, b_{N-1}] & \text{if $w=0$, where $b_i = a_{N-1-i}$} \\
  [b_0, \ldots, b_{w-1}] &  \text{if $w > 0$ and $N\leq w$, where }
      b_i = \begin{cases}
                a_{w-1-i} & \text{if $i \geq w-N$} \\
                0   & \text{if $i \leq w-1-N$}\\
            \end{cases}\\
  \errorX & \text{if $w > 0$ and $N > w$.}
\end{cases}
$$

\noindent In the vacuous case $N=0$ we have $\itobsLE(0,0) = \itobsBE(0,0) = []$, the empty bytestring.

\note{Bytestring to integer conversion.}
\label{note:bstoi}
The \texttt{byteStringToInteger} function converts bytestrings to non-negative
integers.  It takes two arguments:
\begin{itemize}
\item An endianness flag $e$.
\item The bytestring $s$ to be converted.
\end{itemize}
\noindent  
If $e$ is \texttt{True} then the conversion is big-endian, otherwise it is
little-endian.  In both cases the empty bytestring is converted to the integer 0.

\subsubsection{BLS12-381 built-in types}
\label{sec:bls-types-4}

\noindent Supporting the BLS12-381 operations involves adding three new types
and seventeen new built-in functions.  The description of the semantics of these
types and functions is quite complex and requires a considerable amount of
notation, most of which is used only in Sections~\ref{sec:bls-types-4} and~\ref{sec:bls-builtins-4}.

\bigskip
\noindent Table~\ref{table:built-in-types-4} describes three new built-in
types.

\newcommand{\TyMlResult}{\ty{bls12\_381\_mlresult}}
\newcommand{\MlResultDenotation}{H}
\newcommand{\Fq}{\mathbb{F}_q}
\newcommand{\Fqq}{\mathbb{F}_{q^2}}
\newcommand{\FF}{\mathbb{F}_{q^{12}}}

\begin{table}[H]
  \centering
    \begin{tabular}{|l|p{2cm}|l|}
        \hline
        Type & Denotation & Concrete Syntax\\
        \hline
        $\ty{bls12\_381\_G1\_element}$ &   $G_1$ & \texttt{0x[0-9A-Fa-f]\{96\}} \text{(see Note~\ref{note:bls-syntax})}\\
        $\ty{bls12\_381\_G2\_element}$ &   $G_2$ & \texttt{0x[0-9A-Fa-f]\{192\}} \text{(see Note~\ref{note:bls-syntax})}\\
        $\TyMlResult$    &   $\MlResultDenotation$  &  None (see Note~\ref{note:bls-syntax})\\
        \hline
    \end{tabular}
    \caption{Atomic Types}
    \label{table:built-in-types-4}
\end{table}

%% \paragraph{$G_1$ and $G_2$}.
\noindent Here $G_1$ and  $G_2$ are both additive cyclic groups of prime order $r$, where 
$$
r = \mathtt{0x73eda753299d7d483339d80809a1d80553bda402fffe5bfeffffffff00000001}.
$$
        
\paragraph{The fields $\Fq$ and $\Fqq$.}
\noindent To define the groups $G_1$ and $G_2$ we need the finite field $\Fq$ where
\begin{align*}
q = \mathtt{0x}&\mathtt{1a0111ea397fe69a4b1ba7b6434bacd764774b84f38512bf}\\
              &\mathtt{6730d2a0f6b0f6241eabfffeb153ffffb9feffffffffaaab}
\end{align*}

\noindent which is a 381-bit prime. The field $\Fq$ is isomorphic to $\Z_q$,
the ring of integers modulo $q$, and hence there is a natural epimorphism from
$\Z$ to $\Fq$ which we denote by $n \mapsto \bar{n}$.  Given $x \in \Fq$, we
denote by $\tilde{x}$ the smallest non-negative integer $n$ with $\bar{n} = x$.
We sometimes use literal integers to represent elements of $\Fq$ in the obvious
way.

We also make use of the field $\Fqq = \Fq[X]/(X^2+1)$; we may regard $\Fqq$ as
the set $\{a+bu: a,b \in \Fq\}$ where $u^2=-1$.

It is convenient to say that an element $a$ of $\Fq$ is \textit{larger} than
another element $b$ (and $b$ is \textit{smaller} than $a$) if $\tilde{a}
> \tilde{b}$ in $\Z$.  We extend this terminology to $\Fqq$ by saying that
$a+bu$ is larger than $c+du$ if either $b$ is larger than $d$ or $b=d$ and $a$
is larger than $c$.


\paragraph{The groups $G_1$ and $G_2$.}
\noindent There are elliptic curves $E_1$ defined over $\Fq$:
$$
E_1: Y^2 = X^3 + 4
$$

\noindent and $E_2$ defined over $\Fqq$:
$$
E_2: Y^2 = X^3 + 4(u+1).
$$

\noindent $E_1(\Fq)$ and  $E_2(\Fqq)$  are abelian groups under the
usual elliptic curve addition operations as described
in~\cite[III.2]{Silverman-Arithmetic-EC} or~\cite[2.1]{Costello-pairings}.
$G_1$ is a subgroup of $E_1(\Fq)$ and $G_2$ is a subgroup of $E_2(\Fqq)$;
explicit generators for $G_1$ and $G_2$ are given
in~\cite[4.2.1]{IETF-pairing-friendly-curves}.  We denote the identity element
(the point at infinity) in $G_1$ by $\mathcal{O}_{G_1}$ and that in $G_2$ by
$\mathcal{O}_{G_2}$.  Given an integer $n$ and a group element $a$ in $G_1$ or
$G_2$, the scalar multiple $na$ is defined as usual to be $a + \cdots + a$ ($n$
times) if $n>0$ and $-a + \cdots + -a$ ($-n$ times) if $n<0$; $0a$ is the
identity element of the group.

\paragraph{The \texttt{bls12\_381\_MlResult} type.}
\noindent Values of the \texttt{bls12\_381\_MlResult} type are completely
opaque and can only be obtained as a result of \texttt{bls12\_381\_millerLoop}
or by multiplying two existing elements of type \texttt{bls12\_381\_MlResult}.
We provide neither a serialisation format nor a concrete syntax for values of
this type: they exist only ephemerally during computation.  We do not specify
$\MlResultDenotation$, the denotation of $\TyMlResult$, precisely, but it
must be a multiplicative abelian group. See Note~\ref{note:pairing} for more on
this.

\subsubsection{BLS12-381 built-in functions}
\label{sec:bls-builtins-4}

\newcommand{\hash}{\mathsf{hash}}
\newcommand{\compress}{\mathsf{compress}}
\newcommand{\uncompress}{\mathsf{uncompress}}

\setlength{\LTleft}{-4mm} % Shift the table left a bit to centre it on the page
\begin{longtable}[H]{|l|p{5cm}|p{25mm}|c|c|}
    \hline \text{Function} & \text{Signature} & \text{Denotation} & \text{Can}
    & \text{Note} \\ & & & fail?
    & \\ \hline \endfirsthead \hline \text{Function} & \text{Type}
    & \text{Denotation} & \text{Can} & \text{Note}\\ & & & fail?
    & \\ \hline \endhead \hline \caption{BLS12-381 built-in Functions}
    % This caption goes on every page of the table except the last.  Ideally it
    % would appear only on the first page and all the rest would say
    % (continued). Unfortunately it doesn't seem to be easy to do that in a
    % longtable.
    \endfoot
%%    \caption[]{Built-in Functions}
    \caption[]{BLS12-381 built-in Functions (continued)}
    \label{table:built-in-functions-4}
    \endlastfoot
%% G1
    \TT{bls12\_381\_G1\_add}  &
    $[ \ty{bls12\_381\_G1\_element}$,
      \text{\; $\ty{bls12\_381\_G1\_element} ]$}
      \text{\: $ \to \ty{bls12\_381\_G1\_element}$} & $(a,b) \mapsto a+b$  &  No & \\
    \TT{bls12\_381\_G1\_neg}  &
      $ [ \ty{bls12\_381\_G1\_element} ]$  \text{\;\; $\to \ty{bls12\_381\_G1\_element}$} & $a \mapsto -a$  & No & \\
    \TT{bls12\_381\_G1\_scalarMul}  &
    $[ \ty{integer}$,
      \text{\; $\ty{bls12\_381\_G1\_element} ]$}
      \text{\: $ \to \ty{bls12\_381\_G1\_element}$} & $(n,a) \mapsto na$ &  No & \\
    \TT{bls12\_381\_G1\_equal}  &
    $[ \ty{bls12\_381\_G1\_element}$,
      \text{\; $\ty{bls12\_381\_G1\_element} ]$}
      \text{\: $ \to \ty{bool}$} & $=$ &  No & \\
    \TT{bls12\_381\_G1\_hashToGroup}  &
    $[ \ty{bytestring}, \ty{bytestring}]$
      \text{\: $ \to \ty{bls12\_381\_G1\_element}$} & $\hash_{G_1}$ &  Yes & \ref{note:hashing-into-group}\\
    \TT{bls12\_381\_G1\_compress}  &
    $[\ty{bls12\_381\_G1\_element}]$
      \text{\: $ \to \ty{bytestring}$} & $\compress_{G_1}$  &  No & \ref{note:group-compression}\\
    \TT{bls12\_381\_G1\_uncompress}  &
    $[ \ty{bytestring}]$
      \text{\: $ \to \ty{bls12\_381\_G1\_element}$} & $\uncompress_{G_1}$  &  Yes & \ref{note:group-uncompression}\\
    \hline 
%% G2
    \TT{bls12\_381\_G2\_add}  &
    $[ \ty{bls12\_381\_G2\_element}$,
      \text{\; $\ty{bls12\_381\_G2\_element} ]$}
      \text{\: $ \to \ty{bls12\_381\_G2\_element}$} & $(a,b) \mapsto a+b$ &  No & \\
    \TT{bls12\_381\_G2\_neg}  &
      $ [ \ty{bls12\_381\_G2\_element} ]$  \text{\;\; $\to \ty{bls12\_381\_G2\_element}$} & $a \mapsto -a$  & No & \\
    \TT{bls12\_381\_G2\_scalarMul}  &
    $[ \ty{integer}$,
      \text{\; $\ty{bls12\_381\_G2\_element} ]$}
      \text{\: $ \to \ty{bls12\_381\_G2\_element}$} & $(n,a) \mapsto na$ &  No & \\
    \TT{bls12\_381\_G2\_equal}  &
    $[ \ty{bls12\_381\_G2\_element}$,
      \text{\; $\ty{bls12\_381\_G2\_element} ]$}
      \text{\: $ \to \ty{bool}$} & $=$ &  No & \\
    \TT{bls12\_381\_G2\_hashToGroup}  &
    $[ \ty{bytestring}, \ty{bytestring}]$
      \text{\: $ \to \ty{bls12\_381\_G2\_element}$} & $\hash_{G_2}$  &  Yes & \ref{note:hashing-into-group}\\
    \TT{bls12\_381\_G2\_compress}  &
    $[\ty{bls12\_381\_G2\_element}]$
      \text{\: $ \to \ty{bytestring}$} & $\compress_{G_2}$  &  No & \ref{note:group-compression}\\
    \TT{bls12\_381\_G2\_uncompress}  &
    $[ \ty{bytestring}]$
      \text{\: $ \to \ty{bls12\_381\_G2\_element}$} & $\uncompress_{G_2}$  &  Yes & \ref{note:group-uncompression}\\
    \hline 
    \TT{bls12\_381\_millerLoop}  &
    $[ \ty{bls12\_381\_G1\_element}$,
      \text{\; $\ty{bls12\_381\_G2\_element} ]$}
    \text{\: $ \to \TyMlResult$} & $e$ &  No & \ref{note:pairing}\\
    \TT{bls12\_381\_mulMlResult}  &
    $[ \TyMlResult$,
    \text{\; $\TyMlResult]$}
    \text{\: $\to \TyMlResult$} & $(a,b) \mapsto ab$ & No & \ref{note:pairing}\\
    \TT{bls12\_381\_finalVerify}  &
    $[ \TyMlResult$,
    \text{\; $\TyMlResult] \to \ty{bool}$} & $\phi$ & No & \ref{note:pairing}\\
    \hline
\end{longtable}


\note{Hashing into $G_1$ and $G_2$.}
\label{note:hashing-into-group}
The denotations $\hash_{G_1}$ and $\hash_{G_2}$
of \texttt{bls12\_381\_G1\_hashToGroup} and
\texttt{bls12\_381\_G2\_hashToGroup} both take an arbitrary bytestring $b$ (the
\textit{message}) and a (possibly empty) bytestring of length at most 255 known as a \textit{domain
separation tag} (DST)~\cite[2.2.5]{IETF-hash-to-curve} and hash them to obtain a
point in $G_1$ or $G_2$ respectively.  The details of the hashing process are
described in~\cite{IETF-hash-to-curve} (see specifically Section 8.8), except
that
\textbf{we do not support DSTs of length greater than 255}: an attempt to use a
longer DST directly will cause an error.  If a longer DST is required then it
should be hashed to obtain a short DST as described
in~\cite[5.3.3]{IETF-hash-to-curve}, and then this should be supplied as the
second argument to the appropriate \texttt{hashToGroup} function.

%% Some hashing
%% implementations also allow a third argument (an ``augmentation string''), but we
%% do not support this since the same effect can be obtained by appending
%% (prepending?) the augmentation string to the message before hashing.

\newcommand{\ymin}{y_{\text{min}}}
\newcommand{\ymax}{y_{\text{max}}}

\note{Compression for elements of $G_1$ and $G_2$.} 
\label{note:group-compression}
Points in $G_1$ and $G_2$ are encoded as bytestrings in a ``compressed'' format
where only the $x$-coordinate of a point is encoded and some metadata is used to
indicate which of two possible $y$-coordinates the point has.  The encoding
format is based on the Zcash encoding for BLS12-381 points:
see~\cite{Zcash-serialisation} or~\cite[``Serialization'']{BLST-library}
or~\cite[Appendices C and D]{IETF-pairing-friendly-curves}.  In detail,

\begin{itemize}

\item Given an element $x$ of $\Fq$, $\tilde{x}$ can be written as a 381-bit
binary number: $\tilde{x} = \sum_{i=0}^{380}b_i2^i$ with $b_i \in \{0,1\}$.  We
define $\mathsf{bits}(x)$ to be the bitstring $b_{380}\cdots b_0$.

\item A non-identity element of $G_1$ can be written in the form $(x,y)$ with $x,y\in\Fq$.
Not every element $x$ of $\Fq$ is the $x$-coordinate of a point in $G_1$, but
those which are in fact occur as the $x$-coordinate of two distinct points in
$G_1$ whose $y$-coordinates are the negatives of each other.  A similar
statement is true for $\Fqq$ and $G_2$.  In both cases we denote the smaller of
the possible $y$-coordinates by $\ymin(x)$ and the larger by $\ymax(x)$.

\item For $(x,y) \in G_1\backslash \mathcal{O}_{G_1}$ we define
$$
\compress_{G_1} (x,y) = \begin{cases}
\mathsf{1}\mathsf{0}\mathsf{0}\cdot\mathsf{bits}(x) & \text{if $y=\ymin(x)$}\\
\mathsf{1}\mathsf{0}\mathsf{1}\cdot\mathsf{bits}(x) & \text{if $y=\ymax(x)$}\\
\end{cases}
$$
\item We encode the identity element of $G_1$ using
$$
\compress_{G_1}(\mathcal{O}_{G_1}) = \mathsf{1}\mathsf{1}\mathsf{0}\cdot\mathsf{0}^{381},
$$
\noindent where $\mathsf{0}^{381}$ denotes a string of 381 $\mathsf{0}$ bits.
\end{itemize}
\noindent Thus in all cases the encoding of an element of $G_1$ requires exactly 384 bits,
or 48 bytes.

\medskip 

\noindent 
\begin{itemize}
\item Similarly, every non-identity element of $G_2$ can be written
in the form $(x,y)$ with $x,y \in \Fqq$.  We define

$$
\compress_{G_2} (a+bu,y) = \begin{cases}
\mathsf{1}\mathsf{0}\mathsf{0}\cdot\mathsf{bits}(b)\cdot\mathsf{0}\mathsf{0}\mathsf{0}\cdot\mathsf{bits}(a)
& \text{if $y=\ymin(a+bu)$}\\
\mathsf{1}\mathsf{0}\mathsf{1}\cdot\mathsf{bits}(b)\cdot\mathsf{0}\mathsf{0}\mathsf{0}\cdot\mathsf{bits}(a) &
 \text{if $y=\ymax(a+bu)$}\\
\end{cases}
$$

\item The identity element of $G_2$ is encoded using
$$
\compress_{G_2}(\mathcal{O}_{G_2}) = \mathsf{1}\mathsf{1}\mathsf{0}\cdot\mathsf{0}^{765}.
$$

\end{itemize}
\noindent The encoding of an element of $G_2$ requires exactly 768 bits, or 96 bytes.

Note that in all cases the most significant bit of a compressed point is 1.  In
the Zcash serialisation scheme this indicates that the point is compressed;
Zcash also supports a serialisation format where both the $x$- and
$y$-coordinates of a point are encoded, and in that case the leading bit of the
encoded point is 0.  We do not support this format.

\note{Uncompression for elements of $G_1$ and $G_2$.} 
\label{note:group-uncompression}
There are two (partial) ``uncompression'' functions $\uncompress_{G_1}$ and
$\uncompress_{G_2}$ which convert bytestrings into group elements; these are
obtained by inverting the process described in
Note~\ref{note:group-compression}.

\paragraph{Uncompression for $G_1$ elements.}  Given a bytestring $b$, it is checked that
$b$ contains exactly 48 bytes.  If not, then $\uncompress_{G_1}(b) = \errorX$ (ie,
uncompression fails).  If the length is equal to 48 bytes, write $b$ as a
sequence of bits: $b = b_{383} \cdots b_0$.
\begin{itemize}
\item If $b_{383} \neq 1$, then $\uncompress_{G_1}(b) = \errorX$.
\item If $b_{383} = b_{382} = 1$ then
$\uncompress_{G_1}(b) =
\begin{cases}
\mathcal{O}_{G_1} & \text{if $b_{381} = b_{380} = \cdots = b_0 = 0$}\\
\errorX & \text{otherwise}.
\end{cases}$
\item If $b_{383}=1$ and $b_{382}=0$, let $c=\sum_{i=0}^{380}b_i2^i \in \N$.
\begin{itemize}
\item If $c \geq q$, $\uncompress_{G_1}(b) = \errorX$.
\item Otherwise, let $x = \bar{c} \in \Fq$ and let $z = x^3+4$. If $z$ is not a
square in $\Fq$, then $\uncompress_{G_1}(b) = \errorX$.
\item If $z$ is a square then let
$y=\begin{cases}
\ymin(x) & \text{if $b_{381}=0$}\\
\ymax(x) & \text{if $b_{381}=1$}.
\end{cases}$
\item Then $\uncompress_{G_1}(b) = \begin{cases}
(x,y) & \text{if $(x,y) \in G_1$}\\
\errorX & \text{otherwise}.
\end{cases}$
\end{itemize}
\end{itemize}


\paragraph{Uncompression for $G_2$ elements.}  Given a bytestring $b$, it is checked that
$b$ contains exactly 96 bytes.  If not, then $\uncompress_{G_2}(b) = \errorX$ (ie,
uncompression fails).  If the length is equal to 96 bytes, write $b$ as a
sequence of bits: $b = b_{767} \cdots b_0$.
\begin{itemize}
\item If $b_{767} \neq 1$, then $\uncompress_{G_2}(b) = \errorX$.
\item If $b_{767} = b_{766} = 1$ then $\uncompress_{G_2}(b) =
\begin{cases}
\mathcal{O}_{G_2} & \text{if $b_{765} = b_{764} = \cdots = b_0 = 0$}\\
\errorX & \text{otherwise}.
\end{cases}$
\item If $b_{767}=1$ and $b_{766} = 0$, let $c=\sum_{i=0}^{383}b_i2^i$ and $d=\sum_{i=384}^{764}b_i2^{i-384} \in \N$.
\begin{itemize}
\item If $c \geq q$ or $d \geq q$, $\uncompress_{G_2}(b) = \errorX$.
\item Otherwise, let $x = \bar{c}+\bar{d}u \in \Fqq$ and let $z = x^3+4(u+1)$.
If $z$ is not a square in $\Fqq$, then $\uncompress_{G_2}(b) = \errorX$.
\item If $z$ is a square then let
$y=\begin{cases}
\ymin(x) & \text{if $b_{765}=0$}\\
\ymax(x) & \text{if $b_{765}=1$}.
\end{cases}$
\item Then $\uncompress_{G_2}(b) = \begin{cases}
(x,y) & \text{if $(x,y) \in G_2$}\\
\errorX & \text{otherwise}.
\end{cases}$
\end{itemize}
\end{itemize}


\note{Concrete syntax for BLS12-381 types.}
\label{note:bls-syntax}
Concrete syntax for the $\ty{bls12\_381\_G1\_element}$ and
$\ty{bls12\_381\_G2\_element}$ types is provided via the compression and
decompression functions defined in Notes~\ref{note:group-compression}
and~\ref{note:group-uncompression}.  Specifically, a value of type
$\ty{bls12\_381\_G1\_element}$ is denoted by a term of the form \texttt{(con
bls12\_381\_G1\_element 0x...)} where \texttt{...}  consists of 96 hexadecimal
digits representing the 48-byte compressed form of the relevant point.
Similarly, a value of type $\ty{bls12\_381\_G2\_element}$ is denoted by a term
of the form \texttt{(con bls12\_381\_G2\_element 0x...)}  where \texttt{...}
consists of 192 hexadecimal digits representing the 96-byte compressed form of
the relevant point.  \textbf{This syntax is provided only for use in textual
Plutus Core programs}, for example for experimentation and testing.  We do not
support constants of any of the BLS12-381 types in serialised programs on the
Cardano blockchain: see Section~\ref{sec:flat-constants}.  However, for
$\ty{bls12\_381\_G1\_element}$ and $\ty{bls12\_381\_G2\_element}$ one can use
the appropriate uncompression function on a  bytestring constant at runtime:
for example, instead of
$$
\texttt{(con bls12\_381\_G1\_element 0xa1e9a0...)}
$$
write
$$
\texttt{[(builtin bls12\_381\_G1\_uncompress) (con bytestring \#a1e9a0...)]}.
$$

\noindent
No concrete syntax is provided for values of type
$\ty{bls12\_381\_mlresult}$. It is not possible to parse such values, and they
will appear as \texttt{(con bls12\_381\_mlresult <opaque>)} if output by a
program.


\note{Pairing operations.}
\label{note:pairing}
For efficiency reasons we split the pairing process into two parts:
the \texttt{bls12\_381\_millerLoop} and \texttt{bls12\_381\_finalVerify}
functions.  We assume that we have
\begin{itemize}
\item An intermediate multiplicative abelian group $H$.
\item A function (not necessarily itself a pairing) $e: G_1 \times
G_2 \rightarrow \MlResultDenotation$.
\item A cyclic group $\mu_r$ of order $r$.
\item An epimorphism $\psi: \MlResultDenotation \rightarrow \mu_r$ of groups such
that $\psi \circ e: G_1 \times G_2 \rightarrow \mu_r$ is a (nondegenerate,
bilinear) pairing.
\end{itemize}

\noindent Given these ingredients, we define
\begin{itemize}
\item $\denote{\TyMlResult} = \MlResultDenotation$.
\item $\denote{\mathtt{bls12\_381\_mulMlResult}} =$
the group multiplication operation in $\MlResultDenotation$.
\item $\denote{\mathtt{bls12\_381\_millerLoop}} = e$.
\item $\denote{\mathtt{bls12\_381\_finalVerify}} = \phi$,
where
$$
\phi(a,b) = \begin{cases}
               \mathtt{true} & \text{if $\psi(ab^{-1}) = 1_{\mu_r}$} \\
               \mathtt{false} & \text{otherwise.}
            \end{cases}
$$
\end{itemize}

\medskip
\noindent
We do not mandate specific choices for $\MlResultDenotation, \mu_r, e$, and $\phi$, but a
plausible choice would be
\begin{itemize}
\item $\MlResultDenotation = \units{\FF}$.
\item $e$ is the Miller loop associated with the optimal Ate pairing
for $E_1$ and $E_2$~\cite{Vercauteren}.
\item $\mu_r = \{x \in \units{\FF}: x^r=1\}$, the group of $r$th roots of unity in $\FF$.
(There are $r$ distinct $r$th roots of unity in $\FF$ because the embedding
degree of $E_1$ and $E_2$ with respect to $r$ is 12 (see~\cite[4.1]{Costello-pairings}).)
\item $\psi(x) = x^{\frac{q-1}{r}}$.
\end{itemize}

\noindent The functions \texttt{bls12\_381\_millerLoop} and (especially)
\texttt{bls12\_381\_finalVerify} are expected
to be expensive, so their use should be kept to a minimum.  Fortunately most
current use cases do not require many uses of these functions.

\newpage %% Bad page break between table and caption otherwise.
% I tried resetting the note number from V1-builtins here, but that made
% some hyperlinks wrong.  To get note numbers starting at one in each section, I
% think we have to define a new counter every time.
\newcounter{notenumberE}
\renewcommand{\note}[1]{
  \bigskip
  \refstepcounter{notenumberE}
  \noindent\textbf{Note \thenotenumberE. #1}
}

\newcommand\Xand{\mathsf{and}}
\newcommand\Xor{\mathsf{or}}
\newcommand\Xxor{\mathsf{xor}}
\newcommand{\extzero}[1]{\mathtt{0}^*{\cdot}#1}
\newcommand{\extone}[1]{\mathtt{1}^*{\cdot}#1}

\subsection{Batch 5}
\label{sec:default-builtins-5}
The fifth batch of built-in functions adds support for
\begin{itemize}
\item Logical and bitwise operations on bytestrings (see~\cite{CIP-0122} and~\cite{CIP-0123}).
\item The \texttt{RIPEMD-160} hash function.
\end{itemize}

\noindent In the table below most of the functions involve operating on individual bits.
We will often view byetstrings as bitstrings $b_{n-1}\cdots b_0$ with
$b_i \in \{0,1\}$ (and $n$ necessarily a multiple of 8).  Strictly we should use
the functions $\bitsof$ and $\bytesof$ of Section~\ref{sec:notation-bytestrings}
to convert back and forth between bytestrings and bitstrings, but we elide this
for conciseness and reduce ambiguity by denoting bits by lower-case names like
$b$,$c$ and $d$ and bytes by upper-case names like $B$. We denote the
complement of a bit $b \in \{0,1\}$ by $\bar{b}$, so $\bar{\textsf{0}}
= \textsf{1}$ and $\bar{\textsf{1}} = \textsf{0}$.


\setlength{\LTleft}{-6mm}  % Shift the table left a bit to centre it on the page
\begin{longtable}[H]{|l|p{50mm}|p{40mm}|c|c|}
    \hline
    \text{Function} & \text{Signature} & \text{Denotation} & \text{Can} & \text{Note} \\
    & & & fail? & \\
    \hline
    \endfirsthead
    \hline
    \text{Function} & \text{Type} & \text{Denotation} & \text{Can} & \text{Note}\\
    & & & fail? & \\
    \hline
    \endhead
    \hline
    \caption{Built-in functions, batch 5}
    \endfoot
    \caption[]{Built-in functions, batch 5}
    \label{table:built-in-functions-5}
    \endlastfoot
    \TT{andByteString} & $[\ty{bool}, \ty{bytestring}, \ty{bytestring}] $ \text{$\;\; \to \ty{bytestring}$} & $\Xand$ & No & \ref{note:bitwise-logical-ops} \\[2mm]
    \TT{orByteString} & $[\ty{bool}, \ty{bytestring}, \ty{bytestring}] $ \text{$\;\; \to \ty{bytestring}$} & $\Xor$ & No  & \ref{note:bitwise-logical-ops} \\[2mm]
    \TT{xorByteString} & $[\ty{bool}, \ty{bytestring}, \ty{bytestring}] $ \text{$\;\; \to \ty{bytestring}$} & $\Xxor$ & No & \ref{note:bitwise-logical-ops} \\[2mm]
    \TT{complementByteString} & $[\ty{bytestring}] \to \ty{bytestring}$
                              &  $ b_{n-1}\cdots b_0 \mapsto \bar{b}_{n-1}\cdots\bar{b}_0$  & No & \\[2mm]
    \TT{shiftByteString} & $[\ty{bytestring}, \ty{integer}] $ \text{$\;\; \to \ty{bytestring}$} &  $\mathsf{shift}$ & No & \ref{note:shift} \\[2mm]
    \TT{rotateByteString} & $[\ty{bytestring}, \ty{integer}] $ \text{$\;\; \to \ty{bytestring}$} &  $\mathsf{rotate}$ & No & \ref{note:rotate}\\[2mm]
    \TT{countSetBits} & $[\ty{bytestring}] \to \ty{integer}$
    & $b_{n-1} \cdots b_0 $ \text{$\mapsto \left|\{i: b_i =1\}\right|$} & No & \\[2mm]
    \TT{findFirstSetBit} & $[\ty{bytestring}] \to \ty{integer}$ & $\mathsf{ffs}$ & No & \ref{note:ffs}\\[2mm]
    \TT{readBit} & $[\ty{bytestring}, \ty{integer}] $ \text{$\;\; \to \ty{bytestring}$}
                                        & $(b_{n-1}\cdots b_0, i) $ \text{$\mapsto \begin{cases}
                                        b_i & \text{if $0 \leq i \leq n-1$}\\
                                        \errorX & \text{otherwise}
                                        \end{cases}$}
                                        & Yes & \\[14mm]  %% Odd, but it gives reasonable spacing
    \TT{writeBits} & $[\ty{bytestring}, \listOf{\ty{integer}}, $
        \text{$\;\; \listOf{\ty{bool}}] \to \ty{bytestring}$} & $\mathsf{writeBits}$ & Yes & \ref{note:writebits} \\[2mm]
    \TT{replicateByte} & $[\ty{integer}, \ty{integer}] $ \text{$\;\; \to \ty{bytestring}$}
                                        & $\mathsf{replicateByte}$
                                        & Yes & \ref{note:replicateByte}\\[2mm]
    \hline
    \TT{ripemd\_160} & $[\ty{bytestring}] \to \ty{bytestring}$ & Compute a RIPEMD-160 hash \cite{ripemd-2, ripemd-1} & No & \\
\hline
\end{longtable}

\note{Bitwise logical operations.}
\label{note:bitwise-logical-ops}
The bitwise logical operations $\Xand$, $\Xor$, and $\Xxor$ are defined by
extending the usual single-bit logical operations $\wedge$, $\vee$, and $\oplus$
(respectively) to bytestrings. However, there is a subtlety in how bytestrings of
different lengths are dealt with.

\begin{itemize}
\item If the first argument of one of the bitwise logical operations is $\false$ 
then the longer bytestring is truncated on the left to have the same length as
the shorter one.
\item If the first argument $\true$ 
then the shorter bytestring is extended on the left to have the same length as
the longer one.  In the case of $\Xand$ the shorter bytestring is extended
with \texttt{1} bits and in the case of $\Xor$ and $\Xxor$ it is extended
with \texttt{0} bits.
\end{itemize}

\noindent To specify the extension process mentioned above it is helpful to define two
\textit{infinite} bitstrings $\extzero{s}$ and $\extone{s}$ which extend a finite bitstring $s$
to the left:

$$
(\extzero{s})_i =
\begin{cases}
  s_i & \text{if $0 \leq i \leq n-1$}\\
  \mathtt{0} & \text{if $i \geq n$}
\end{cases}
$$

$$
(\extone{s})_i =
\begin{cases}
  s_i & \text{if $0 \leq i \leq n-1$}\\
  \mathtt{1} & \text{if $i \geq n$}
\end{cases}
$$

\noindent The denotations of the bitwise logical functions are now defined on
bitstring representations of bytestrings by

\begin{align*}
\Xand\;(\false, b_{m-1} \cdots b_0, c_{n-1} \cdots c_0) &= d_{l-1} \cdots d_0 
\quad \text{where $l=\min\{m,n\}$ and $d_i = b_i \wedge c_i$}\\
\Xand\;(\true, a_{m-1} \cdots a_0, b_{n-1} \cdots b_0) &= d_{l-1} \cdots d_0
\quad \text{where $l=\max\{m,n\}$ and $d_i = (\extone{a})_i \wedge(\extone{b})_i$}
\end{align*}


\begin{align*}
\Xor\;(\false, b_{m-1} \cdots b_0, c_{n-1} \cdots c_0) &= d_{l-1} \cdots d_0
  \quad \text{where $l=\min\{m,n\}$ and $d_i = b_i \vee c_i$}\\
\Xor\;(\true, b_{m-1} \cdots b_0, c_{n-1} \cdots c_0) &= d_{l-1} \cdots d_0 
\quad \text{ where $l=\max\{m,n\}$ and $d_i = (\extzero{a})_i \vee (\extzero{b})_i$}
\end{align*}


\begin{align*}
\Xxor\;(\false, b_{m-1} \cdots b_0, c_{n-1} \cdots c_0) &= d_{l-1} \cdots d_0 
\quad \text{where $l=\min\{m,n\}$ and $d_i = b_i \oplus c_i$}\\
\Xxor\;(\true, b_{m-1} \cdots b_0, c_{n-1} \cdots c_0) &= d_{l-1} \cdots d_0 
\quad \text{where $l=\max\{m,n\}$ and $d_i = (\extzero{a})_i \oplus (\extzero{b})_i$.}
\end{align*}


\note{Shifting bits.}
\label{note:shift}
The \texttt{shiftByteString} builtin takes a bytestring $s$ and an integer $k$ and
shifts the bits of the bytestring $k$ places to the left if $k \geq 0$ and $k$
places to the right if $k < 0$, replacing vacated bits with \texttt{0}.  The
length of the output bytestring is the same as that of the input.  The
denotation (defined on the bitstring representation of $s$) is

$$
\mathsf{shift}\;(b_{n-1} \cdots b_0, k) =
  d_{n-1} \cdots d_0 \quad \text{where }
  d_i = \begin{cases}
     b_{i-k} & \text{if $0 \leq i-k \leq n-1$ }\\
     \texttt{0} & \text{otherwise}\\
\end{cases}
$$

\note{Rotating bits.}
\label{note:rotate}
The \texttt{rotateByteString} builtin takes a bytestring $s$ and an integer $k$
and rotates the bits of $s$ $k$ places to the left if $k \geq 0$ and $k$ places
to the right if $k < 0$.  The length of the output bytestring is the same as
that of the input.  The denotation of
\texttt{rotateByteString} (defined on the bitstring representation of $s$) is 

$$
\mathsf{rotate}\;(b_{n-1} \cdots b_0, k) = d_{n-1}\cdots d_0 \quad\text{where $d_i = b_{(i-k)\bmod n}$}
$$

\note{Finding the first set bit in a bytestring.}
\label{note:ffs}
The \texttt{findFirstSetBit} builtin returns the index of the first nonzero bit
in a bytestring, counting from the \textit{right}. If the bytestring consists
entirely of zeros then it returns $-1$.  The denotation
$\mathsf{ffs}: \b^* \rightarrow \Z$ is

$$
\mathsf{ffs}(b_{n-1}\cdots b_0) =
\begin{cases}
  -1 & \text{if $b_i = \mathtt{0}$ for $0 \leq i \leq n-1$}\\
  \min{\{i: b_i \ne \mathtt{0}\}} & \text{otherwise}
  \end{cases}
$$


\note{\textsf{Writing bits.}}
\label{note:writebits}
The denotation $\mathsf{writeBits}$ of the \texttt{writeBits} builtin takes a
bytestring $b$, a list of $J$ of integer indices, and a list $V$ of booleans.
If the lists are of different lengths then the longer one is truncated to make
them both the same length.  After this, an error occurs if any of the indices in
$J$ is not a valid bit index for $b$.  If all of the indices are within bounds
then for each index $j$ in $J$ the $j$-th bit of the bytestring is updated
according to the value of the corresponding element of $V$ (\texttt{0}
for \textsf{false}, \texttt{1} for \textsf{true}).  Conceptually the updates
take place in sequence according to their positions in the lists: if a bit at a
particular position is updated multiple times then the final update takes
precedence.

\smallskip
\noindent Formally, 
$$
\mathsf{writeBits}(s, [j_0 , \ldots, j_{k-1}], [v_0, \ldots, v_{l-1}]) =
    \mathsf{writeBits}^{\prime} (s, [(j_0, v_0) , \ldots, (j_{\min\{k,l\}-1}, v_{\min\{k,l\}-1})])
$$

$$
\mathsf{writeBits}^{\prime}(b_{n-1}{\cdots}b_0, [(j_0,v_0), \ldots, (j_{l-1}, v_{l-1})]) = 
\begin{cases}
\errorX & \text{if $\exists k \in \{0, \ldots, l-1\}$ such that $j_k<0$ or $j_k \geq n$} \\
d_{n-1}{\cdots}d_0 & \text{otherwise}
\end{cases}
$$

\noindent where
$$
  d_i =
  \begin{cases}
    b_i & \text{if $i \not{\in} \{j_0, \ldots, j_{l-1}\}$}\\
    w & \text{otherwise, where $m = \max\{k : j_k = i\}$ and }
      w=\begin{cases}
      \mathtt{0} & \text{if $v_m = \false$}\\
      \mathtt{1} & \text{if $v_m = \true$}.
    \end{cases}
  \end{cases}
$$


\note{Replicating bytes.}
\label{note:replicateByte}
The \texttt{replicateByte} builtin takes a length $l$ between 0 and 8192 and an
integer $B$ between 0 and 255 and produces a bytestring of length $l$.  It fails
if either argument is outside the required bounds.  The denotation is

$$
\mathsf{replicateByte}(l,B) =
\begin{cases}
B_0{\cdots}B_{l-1} \quad\text{ ($B_i=B$ for all $i$}) & \text{if $0 \leq l \leq 8192$ and $B \in \B$}\\
\errorX      & \text{otherwise}.
\end{cases}
$$
\noindent Note that unlike the other denotations in this section we are
viewing the output as a bytestring here, not a bitstring.

% I tried resetting the note number from V1-builtins here, but that made
% some hyperlinks wrong.  To get note numbers starting at one in each section, I
% think we have to define a new counter every time.
\newcounter{notenumberF}
\renewcommand{\note}[1]{
  \bigskip
  \refstepcounter{notenumberF}
  \noindent\textbf{Note \thenotenumberF. #1}
}
\newpage
\subsection{Batch 6}
\label{sec:default-builtins-6}
The sixth batch of built-in types and functions provides
\begin{itemize}
\item A modular exponentiation function (see CIP-0109~\cite{CIP-0109}).
\item A function to drop a specified number of items from the front of a built-in list (see CIP-0132~\cite{CIP-0132}).
\item A built-in array type and functions for working with arrays (see CIP-0138)~\cite{CIP-0138})
\item Two functions to provide efficient multi-scalar multiplication for the BLS12-381 elliptic curves (see CIP-0133~\cite{CIP-0133}).
\item A built-in \texttt{value} type to represent quantities of currencies on the
      Cardano blockchain, together with a number of operations on values. These
      are described separately from the rest of the functions, in
      Sections~\ref{sec:built-in-value-type} and \ref{sec:built-in-value-functions}. See CIP-0153~\cite{CIP-0153}.
\end{itemize}

\subsubsection{The \texttt{array} type operator}
\label{sec:built-in-type-operators-6}
The \texttt{array} type operator is described in Table~\ref{table:array-type-operator}.

\begin{table}[H]
  \centering
    \begin{tabular}{|l|p{14mm}|l|l|}
        \hline
        Operator $\op$ & $\left|\op\right|$  & Denotation & Concrete Syntax\\
        \hline
        \texttt{array}
          & 1
          & $\denote{\arrayOf{t}} = \denote{t}^*$
          & See below\\
        \hline
        \end{tabular}
   \caption{Array type operator, Batch 6}
    \label{table:array-type-operator}
\end{table}

\paragraph{Arrays.}
Arrays are finite sequences of built-in values of a particular type with the
expectation that looking up the $n$th element takes constant time (unlike
lists).

An array of type $\texttt{array}(t)$ is written as a syntactic list
\texttt{[$c_1, \ldots, c_n$]} where each $c_i$ lies in $\bitc_t$.
Some valid constant expressions are thus:
\begin{verbatim}
   (con (array bool) [True, False, True])
   (con (array integer) [])
\end{verbatim}

\subsubsection{Built-in functions (1)}
\label{sec:built-in-functions-6-1}
The non-\texttt{value}-related built-in functions in Batch 6 are defined in
Table~\ref{table:built-in-functions-6-1}.

\setlength{\LTleft}{-23mm}  % Shift the table left a bit to centre it on the page
\renewcommand*{\arraystretch}{1.25}  %% Stretch the space between the rows by a factor of 25%
\begin{longtable}[H]{|l|p{50mm}|p{58mm}|c|c|}
    \hline
    \text{Function} & \text{Signature} & \text{Denotation} & \text{Can} & \text{Note} \\
    & & & fail? & \\
    \hline
    \endfirsthead
    \hline
    \text{Function} & \text{Type} & \text{Denotation} & \text{Can} & \text{Note}\\
    & & & fail? & \\
        \hline
    \endhead
    \hline
    \caption{Built-in functions, Batch 6 (1)}
    \endfoot
    \caption[]{Built-in functions, Batch 6 (1)}
    \label{table:built-in-functions-6-1}
    \endlastfoot
    \TT{expModInteger}        & $[\ty{integer}, \ty{integer}, \ty{integer}]$ \text{$\;\;\; \to \ty{integer}$}
        & $\mathsf{expMod} $  & Yes & \ref{note:exp-mod-integer}\\
    \TT{dropList}        & $[\forall a_\#, \ty{integer}, \listOf{a_\#}]$ \text{$\;\;\; \to \listOf{a_\#}$}
        & $(k,[x_1,\ldots,x_n])$
        \smallskip
        \newline
        \text{$\;\;\mapsto \left\{ \begin{array}{ll}
            [x_1,\ldots, x_n]      &  \text{if $k \leq 0$} \\ \relax %
            % Without \relax the [ at the start of the next line is regarded as an argument to \\
            [x_{k+1}, \ldots, x_n]  & \text{if $1  \leq k \leq n-1$} \\ \relax
            []                     &\text{if $k \geq n$}\\
        \end{array}\right.$}
        & No & \\
    \TT{lengthOfArray}
      & $[\forall a_\#, \arrayOf{a_\#}] \to \ty{integer}$
      & $[] \mapsto 0$
        \newline
        $[x_1,\ldots,x_n] \mapsto n\ (n \geq 1)$
      & No
      & \ref{note:array-semantics}\\
    \TT{listToArray}
      & $[\forall a_\#, \listOf{a_\#}] \to \arrayOf{a_\#}$
      & $[] \mapsto []$
        \newline
        $[x_1,\ldots,x_n] \mapsto [x_1,\ldots,x_n]\ (n \geq 1)$
      & No
      & \ref{note:array-semantics}\\
    \TT{indexArray}
      & $[\forall a_\#, \arrayOf{a_\#}, \ty{integer}]$ \text{$\;\;\; \to \ty{a_\#}$}
      &  $([x_1,\ldots,x_n], k)$
        \smallskip
        \newline
        \text{$\;\;\mapsto \left\{ \begin{array}{ll}
            \errorX   & \text{if $k < 0$} \\ \relax %
            x_{k+1}   & \text{if $0 \leq k \leq n-1$} \\ \relax
            \errorX   & \text{if $k > n-1$}\\
        \end{array}\right.$}
      & Yes & \ref{note:array-semantics}\\[2mm]
    \TT{bls12\_381\_G1\_multiScalarMul}  &
    $[ \listOf{\ty{integer}}$,
      \text{\; ${\listOf{\ty{bls12\_381\_G1\_element}}} ]$}
      \text{\: $ \to \ty{bls12\_381\_G1\_element}$}
      & $([n_1, \ldots, n_k],[a_1, \ldots, a_l])$
      \newline
      \text{$\: \mapsto n_1a_1+\cdots+n_ra_r$ \ where $r=\min(k,l)$}
      & No & \ref{note:msm-semantics}\\[2mm]
    \TT{bls12\_381\_G2\_multiScalarMul}  &
    $[ \listOf{\ty{integer}}$,
      \text{\; ${\listOf{\ty{bls12\_381\_G2\_element}}} ]$}
      \text{\;\; $ \to \ty{bls12\_381\_G2\_element}$}
      & $([n_1, \ldots, n_k],[a_1, \ldots, a_l])$
      \newline
      \text{$\: \mapsto n_1a_1+\cdots+n_ra_r$ \ where $r=\min(k,l)$}
      & No & \ref{note:msm-semantics}\\[2mm]
\hline
\end{longtable}

\note{Modular Exponentiation.}
\label{note:exp-mod-integer}
The \texttt{expModInteger} function performs modular exponentiation.  Its denotation
$\mathsf{expMod}: \Z \times \Z \times \Z \rightarrow \Z_{\errorX}$ is given by

$$
\mathsf{expMod}(a,e,m) =
  \begin{cases}
     \modfn(a^e,m) & \text{if $m > 1$ and $e \geq 0$}\\
     \min\{r \in \N: ra^{-e} \equiv 1\pmod{m}\}\  & \text{if $m > 1, e < 0$ and $\mathrm{gcd}(a,m) = 1$}\\
     \errorX & \text{if $m>1,e<0$, and $\mathrm{gcd}(a,m) \neq 1$}\\
     0 & \text{if $m=1$}\\
     \errorX & \text{if $m \leq 0$.}
  \end{cases}
$$

\noindent The fact that this is well defined for the case $e<0$ and $\mathrm{gcd}(a,m) = 1$
follows, for example, from Proposition I.3.1 of~\cite{Koblitz-GTM}.  See
Note~\ref{note:integer-division-functions} of
Section~\ref{sec:built-in-functions-1} for the definition of $\modfn$.

\note{Semantics of arrays.}
\label{note:array-semantics}
The denotation of the array type is identical to that of lists, but we expect
that in practice arrays will be implemented in such a way that the
\texttt{lengthOfArray} and \texttt{indexArray} functions will be constant time and
\texttt{listToArray} may be nontrivial, requiring linear time.

\note{BLS12-381 multi-scalar multiplication.}
\label{note:msm-semantics}
We have added two multi-scalar-multiplication operations for BLS12-381
types. Given lists of scalars $[n_1, \ldots, n_k]$ and group elements
$[a_1, \ldots, a_l]$ these compute the group sum of the scalar products $n_ia_i$
with $1 \leq i \leq \min(k,l)$. The input lists are allowed to be of different
lengths but the longer one is effectively truncated to make them the same
length.  If either list is empty then the result of the function is the zero
element of the relevant group.  Multi-scalar multiplication can already be
performed using the standard group addition and scalar multiplication functions
from Batch 4, but having dedicated functions allows for significantly more
efficient implementation.

\newcommand{\BB}{\mathtt{B}\;}
\newcommand{\II}{\mathtt{I}\;}
\newcommand{\flatten}[1]{\ensuremath{#1}\text{\textdownarrow}}
\newcommand{\unflatten}[1]{\ensuremath{#1}\text{\textuparrow}}
% If you use \downarrow and \uparrow there's too much space before them.

\subsubsection{The built-in \texttt{value} type}
\label{sec:built-in-value-type}

\paragraph{Built-in \textsf{Values}.}
The built-in \texttt{value} type represents the \textsf{Value} type used to
represent quantities of currencies in Cardano's EUTXO model: see~\cite{EUTXOma};
we provide a built-in type to represent these and a number of operations for
working with them, as proposed in CIP-0153~\cite{CIP-0153}.

A \texttt{value} involves a finite set of \textit{currency identifiers}, each of
which has a finite number of associated \textit{token identifiers}, each of
which has an associated integral \textit{quantity}. Various other names are used
for these: for example the Cardano ledger and the Aiken language call currency
identifiers \textit{policy IDs} and token identifiers \textit{asset names}, and
the Plinth and Plutarch languages call currency identifiers \textit{currency
symbols} and token identifiers \textit{token names}.

Currency and token identifiers are bytestrings, and in a
particular \texttt{value} a given token identifier may be associated with
multiple currency identifiers.  We model \texttt{value}s as finitely supported
maps from (currency identifier, token identifier) pairs to integers: see
Table~\ref{table:built-in-value-type} (and recall that $\B^*$ denotes the set of
all bytestrings, ie finite sequences of bytes).  For brevity we will call such
finitely supported maps \textit{values} in this section.  In
practice values will be implemented as some finite data structure in
which only a finite number of currency and token identifiers appear, but for
specification purposes it is convenient to allow a value to
contain \textit{all} identifiers, but with only a finite number of quantities
nonzero.  In practice we only allow values $v$ such that for $c, t \in \B^*$
with $v(c,t) \neq 0$, it is the case that $\length(c) \leq 32$, $\length(t) \leq
32$ and $-2^{127} \leq q \leq 2^{127}-1$.  We will call such
values \textit{small values}; none of the built-in functions we provide allow a
non-small value to be created, and non-small values will also cause a failure
during deserialisation from the \texttt{flat} format (see
Appendix~\ref{appendix:flat-serialisation}) and during parsing of textual UPLC.


\begin{table}[H]
  \centering
    \begin{tabular}{|l|p{3cm}|l|}
        \hline
        Type & Denotation & Concrete Syntax\\
        \hline
        $\ty{value}$ & $\mathsf{FinSup}(\B^* \times \B^*, \Z)$ & See below\\
        \hline
    \end{tabular}
    \caption{Atomic types, Batch 6}
    \label{table:built-in-value-type}
\end{table}

\noindent To obtain a syntactic representation of the abstract value type
we define a concrete ``flattened'' form which represents a value as a
nested association list.  Given $v \in \mathsf{FinSup}(\B^* \times \B^*, \Z)$,
let
$$C_v
= \{c \in \B^*: \text{there exists $t \in \B^*$ with $v(c,t) \neq 0\}$.}
$$
\noindent
Since $v$ is finitely supported this set is finite, and since the set $\B^*$ of
all bytestrings is totally ordered (lexicographically), we can write
$$
 C_v = \{c_1, \ldots, c_n\}\ \text{with $c_1 < c_2 < \cdots < c_n$}
$$
\noindent with $n\geq 0$, the case $n=0$ corresponding to the zero map.
Similarly, for $1 \leq i \leq n$ we let
$$T_v(i) = \{t \in \B^*: v(c_i,t) \neq 0\};
$$
\noindent
each of these sets is nonempty and we can write
$$
   T_v(i) = \{t_{i1}, t_{i2}, \ldots, t_{ik_i}\}\ \text{with $t_{i1} < t_{i2} < \cdots < t_{ik_i}$}
$$
\noindent with $k_i \geq 1$ for $1 \leq i \leq n$. We now define $\flatten{v} \in (\B^* \times (\B^* \times \Z)^*)^*$ by
\begin{align*}
 \flatten{v} =
 [&(c_1, [(t_{11}, v(c_1, t_{11})), (t_{12}, v(c_1, t_{12})), \ldots, (t_{1k_1}, v(c_1, t_{1k_1}))],\\
  &(c_2, [(t_{21}, v(c_2, t_{21})), (t_{22}, v(c_2, t_{22})), \ldots, (t_{2k_2}, v(c_2, t_{2k_2}))],\\
  &\ldots,\\
  &(c_n, [(t_{n1}, v(c_n, t_{n1})), (t_{n2}, v(c_n, t_{n2})), \ldots, (t_{nk_n}, v(c_n, t_{nk_n}))]]
\end{align*}%
\nomenclature[Ya]{$\flatten{v}$}{Concrete representation of a \texttt{value} $v$ using nested association lists}

\noindent
Conversely, given a list
\begin{align*}
 l =  [&(c_1, [(t_{11}, q_{11}), (t_{12}, q_{12}), \ldots, (t_{1k_1}, q_{1k_1})]),\\
       &(c_2, [(t_{21}, q_{21}), (t_{22}, q_{22}), \ldots, (t_{2k_2}, q_{2k_2})]),\\
       &\ldots,\\
       &(c_n, [(t_{n1}, q_{n1}), (t_{n2}, q_{n2}), \ldots, (t_{nk_n}, q_{nk_n})])]
\end{align*}

\noindent in $(\B^* \times (\B^* \times \Z)^*)^*$ with
\begin{itemize}
 \item $n \geq 0$
 \item $c_1 < c_2 < \cdots <c_n$
 \item $k_i \geq 1$ for $1 \leq i \leq n$
 \item $t_{i1} < t_{i2} < \cdots < t_{ik_i}$ for $1 \leq i \leq n$
 \item $q_{ij} \in [-2^{127}, 2^{127}-1]\setminus 0$ for all $i$ and $j$
\end{itemize}
(we call a list satisfying these conditions a \textit{well-formed} list) we can
define a finitely supported map
$\unflatten{l} \in \mathsf{FinSup}(\B^* \times \B^*, \Z)$ by
$$
  \unflatten{l}(c,t) =
  \left\{\begin{array}{ll}
  q_{ij} & \text{if $c=c_i$ for some $i$ and $t=t_{ij}$ for some $i$ and $j$}\\
  0 &\text{otherwise.}
  \end{array}\right.
$$%
\nomenclature[Yb]{$\unflatten{l}$}{A nested list $l$ converted into a \texttt{value}}

\noindent
This function is well defined because the well-formedness requirements ensure
that a given pair $(c,t)$ appears at most once in the nested list.  If $l$ is
not well-formed then we put
$$
 \unflatten{l} = \errorX.
$$

\noindent
We have $\unflatten{\flatten{v}} = v$ for all small values $v$ and
$\flatten{\unflatten{l}} = l$ for all well-formed $l$, the latter because
well-formed lists provide a canonical representation of values: for each such
value $v$ there is exactly one well-formed list $l$ with $\unflatten{l} = v$.
Without the restrictions we could permute the entries of the lists and add
currencies with empty token lists or tokens with zero quantities to obtain
different representations of the same value.

\paragraph{Concrete syntax.}
In the concrete syntax, constants of type \texttt{value} are represented
by strings of the form of the form
\begin{align*}
 \mathtt{[}
  & \text{\texttt{($C_1$, [($T_{11}$, $Q_{11}$), ($T_{12}$, $Q_{12}$), $\ldots$, ($T_{1k_1}$, $Q_{1k_1}$)]),}}\\
  & \text{\texttt{($C_2$, [($T_{21}$, $Q_{21}$), ($T_{22}$, $Q_{22}$), $\ldots$, ($T_{2k_2}$, $Q_{2k_2}$)]),}}\\
  & \text{\texttt{$\ldots$,}} \\
  & \text{\texttt{($C_n$, [($T_{n1}$, $Q_{n1}$), ($T_{n2}$, $Q_{n2}$), $\ldots$, ($T_{nk_n}$, $Q_{nk_n}$)])]}}
\end{align*}

\noindent where each $C_i$ and $T_i$ are literal bytestrings
described by the regular expression \texttt{\#([0-9A-Fa-f][0-9A-Fa-f])*} with at
most 32 pairs of hexadecimal digits and each $Q_{ij}$ is a non-zero literal integer
in the range $[-2^{127}, 2^{127}-1]$. The currency IDs and the token IDs for
each currency must be in strictly ascending order.

\medskip
\noindent
The denotation of the associated \texttt{value} is 
\begin{align*}
 \unflatten{[
  & (\denote{C_1}, [(\denote{T_{11}}, \denote{Q_{11}}), (\denote{T_{12}}, \denote{Q_{12}}),
      \ldots, (\denote{T_{1k_1}}, \denote{Q_{1k_1}})]),\\
  & (\denote{C_2}, [(\denote{T_{21}}, \denote{Q_{21}}), (\denote{T_{22}}, \denote{Q_{22}}),
      \ldots, (\denote{T_{2k_2}}, \denote{Q_{2k_2}})]),\\
  & \ldots, \\
  & (\denote{C_n}, [(\denote{T_{n1}}, \denote{Q_{n1}}), (\denote{T_{n2}}, \denote{Q_{n2}}),
      \ldots, (\denote{T_{nk_n}}, \denote{Q_{nk_n}})])]}
  \in \mathsf{FinSup}(\B^* \times \B^*, \Z)
\end{align*}

\noindent More intuitively (but less precisely), this maps every pair $(C_i, T_{ij})$
occurring in the constant to $Q_{ij}$ and everything else to 0.

\medskip
\noindent
Some examples of valid \texttt{value} constants are
\begin{verbatim}
  (con value []) 
  (con value [(#1234AB, [(#, 2), (#01, 44), (#05, -3)])])
  (con value [(#11, [(#, 2), (#01, 1), (#05, 3)]), 
              (#3333, [(#0f,  4), (#1000, -1)])])
\end{verbatim}

\noindent
The expressions
\begin{verbatim}
  (con value [(#33,[(#1234, 455)]), (#2222, [(#9a, -7)])])
\end{verbatim}
\noindent and
\begin{verbatim}
  (con value [(#2222, [(#9a, 44), (#F0, 0), (#FF, 22)])])
\end{verbatim}
\noindent are not valid, the first because the currency IDs are not in ascending order
and the second because the token with ID \texttt{\#F0} appears with quantity 0.


\subsubsection{Built-in functions operating on \texttt{value}s}
\label{sec:built-in-value-functions}

The built-in functions operating on \texttt{value}s are defined in
Table~\ref{table:built-in-value-functions}.  The inputs to most of these are
restricted in size and to make it easier to specify them we introduce the sets
$\B^*_{32}$ of bytestrings of length at most 32 and $Q$ of 128-bit integers:

$$\B^*_{32} = \{[c_0, \ldots, c_{n-1}] \in \B^*: n \leq 32\}$$

$$Q = \{n \in \Z: -2^{127} \leq n \leq 2^{127}-1\}.$$




\setlength{\LTleft}{-10mm}  % Shift the table left a bit to centre it on the page
\renewcommand*{\arraystretch}{1.25}  %% Stretch the space between the rows by a factor of 25%
\begin{longtable}[H]{|l|p{50mm}|p{58mm}|c|c|}
    \hline
    \text{Function} & \text{Signature} & \text{Denotation} & \text{Can} & \text{Note} \\
    & & & fail? & \\
    \hline
    \endfirsthead
    \hline
    \text{Function} & \text{Type} & \text{Denotation} & \text{Can} & \text{Note}\\
    & & & fail? & \\
        \hline
    \endhead
    \hline
    \caption{Batch 6: built-in \texttt{value} functions}
    \endfoot
    \caption[]{Batch 6: built-in \texttt{value} functions}
    \label{table:built-in-value-functions}
    \endlastfoot
\hline  % Value builtins

  \TT{insertCoin} & $[ \ty{bytestring}, \ty{bytestring}, $
                       \text{$\:  \ty{integer}, \ty{value} ]$} $\to \ty{value}$
                  & $(c, t, n, v)$
                  \smallskip
                  \newline
                  $\mapsto \left\{
                  \begin{array}{ll}
                  v[(c,t) \mapsto 0] & \text{if $n = 0$}\\
                  v[(c,t) \mapsto n] & \text{if $c,t\in B^*_{32}$, $n \in Q$}\\
                  \errorX & \text{otherwise}
                  \end{array}\right.$
                  & Yes & \ref{note:insertCoin-semantics} \\[2mm]
  \TT{lookupCoin} & $[ \ty{bytestring}, \ty{bytestring}, $
                       \text{$\: \ty{value} ]$} $\to \ty{integer}$
                  & $ (c,t,v) \mapsto v(c,t) $
                  & No & \ref{note:lookupCoin-semantics} \\[2mm]
  \TT{unionValue} & $[ \ty{value}, \ty{value} ] \to \ty{value}$
                  & $ \mathsf{unionValue} $
                  & Yes  & \ref{note:unionValue-semantics} \\[2mm]
  \TT{valueContains} & $[ \ty{value}, \ty{value} ] \to \ty{bool}$ & $\mathsf{valueContains}$
                  & Yes & \ref{note:valueContains-semantics} \\[2mm]
  \TT{scaleValue} & $[ \ty{integer}, \ty{value} ] \to \ty{value}$
                  & $\mathsf{scaleValue}$
                  & Yes & \ref{note:scaleValue-semantics} \\[2mm]
  \TT{valueData} & $[ \ty{value} ] \to \ty{data}$ &  \textsf{valueData} & Yes & \ref{note:valueData-semantics} \\[2mm]
  \TT{unValueData} & $[ \ty{data} ] \to \ty{value}$ & \textsf{unValueData} & Yes & \ref{note:unValueData-semantics} \\[2mm]
\hline
\end{longtable}



\note{Semantics of \texttt{insertCoin}.}
\label{note:insertCoin-semantics}
The \texttt{insertCoin} function updates the quantity $q$ associated with a
(currency identifier, token identifier) pair $(c,t)$ in a value $v$. If $q=0$
then the function returns $v$ with the entry for $(c,t)$ (if present) deleted;
if there is no entry for $(c,t)$ then the function succeeds and $v$ is returned
unaltered. When $q=0$ there is no limit on the size of the identifiers and the
function always succeeds.  If $q$ is nonzero and lies in $[-2^{127}, 2^{127}-1]$
and both $c$ and $t$ are no longer than 32 bytes then the function returns $v$
updated to map $(c,t)$ to $q$; any previous quantity associated with $(c,t)$ is
discarded.  If $q \neq0$ and any of $c$, $t$, or $q$ is out of bounds then the
function fails.

\note{Semantics of \texttt{lookupCoin}.}
\label{note:lookupCoin-semantics}
The \texttt{lookupCoin} function always accepts currency and token identifiers
of any length; however there is no way to create a value where an identifier
more than 32 bytes long is associated with a nonzero quantity, so
if \texttt{lookupCoin} is called with such an identifier it will always succeed
and return 0.

\note{Semantics of \texttt{unionValue}.}
\label{note:unionValue-semantics}
The result of \texttt{unionValue} is obtained by adding together corresponding
quantities in its inputs, provided that all of the added quantities are in the
signed 128-bit range $[-2^{127}, 2^{127}-1]$: its denotation
\textsf{unionValue} is defined by
  $$
  \mathsf{unionValue}(v, w) =
  \left\{
  \begin{array}{ll}
        v+w & \text{if $v(c,t)+w(c,t) \in Q$ for all $c$ and $t$}\\
        \errorX & \text{otherwise}\\
  \end{array}\right.
  $$
where $v+w \in \mathsf{FinSup}(\B^* \times \B^*, \Z)$ is defined by $(v+w)(c,t)
= v(c,t)+w(c,t)$ for all $c$ and $t$.

\note{Semantics of \texttt{valueContains}.}
\label{note:valueContains-semantics}
For \texttt{valueContains} $v$ $w$, an error occurs if either $v$ or $w$
contains a negative quantity.  If all quantities are non-negative then the
function returns \texttt{true} if every quantity in $v$ is greater than or equal
to the corresponding quantity in $w$, otherwise it returns \texttt{false}: its
denotation
\textsf{valueContains} is defined by
  $$
  \mathsf{valueContains}(v,w)  =
  \left\{ \begin{array}{ll}
  \textsf{true} & \text{if  $v(c,t) \geq w(c,t) \geq 0$ for all $c$ and $t$}\\
  \textsf{false} & \text{if $v(c_0,t_0) < w(c_0,t_0)$ for some $c_0$ and $t_0$ and $v(c,t), w(c,t) \geq 0$ for all $c$ and $t$}\\
  \errorX & \text{otherwise.}\\
  \end{array}\right.
  $$

\note{Semantics of \texttt{scaleValue}.}
\label{note:scaleValue-semantics}
The result of \texttt{scaleValue} $n$ $v$ is $v$ with all quantities multiplied
by $n$, provided that all of the scaled quantities are in the signed 128-bit
range: its denotation
\textsf{scaleValue} is defined by
  $$
  \mathsf{scaleValue}(n,v) =
  \left\{
  \begin{array}{ll}
        n*v & \text{if $(n*v)(c,t) \in Q$ for all $c$ and $t$}\\
        \errorX & \text{otherwise}\\
  \end{array}\right.
  $$
  where $n*v \in \mathsf{FinSup}(\B^* \times \B^*, \Z)$ is defined by
  $(n*v)(c,t) = n*v(c,t)$ for all $c$ and $t$.

\note{Semantics of \texttt{valueData}.}
\label{note:valueData-semantics}
The \texttt{valueData} function converts a \texttt{value} object to an object of
type \texttt{data} (see Section~\ref{sec:default-builtins-1}).  It returns
a \texttt{Map} (the \textit{outer map}) containing a list of pairs, the first
element of each pair being a \texttt{B} object containing a currency identifier,
and the second element (an \textit{inner map}) containing a list of (token
identifier, quantity) pairs.  The elements of the outer and inner maps are
sorted lexicograpically by identifier, there are no repeated identifiers in any
map, and the quantities are all nonzero.  The function fails if there are more
than 40,000 token IDs in total in the input \texttt{value}.

\medskip
\noindent Formally, given a value $v$ with

\begin{align*}
 \flatten{v}=
 [&(c_1, [(t_{11}, q_{11}), (t_{12}, q_{12}), \ldots, (t_{1k_1}, q_{1k_1})]),\\
  &(c_2, [(t_{21}, q_{21}), (t_{22}, q_{22}), \ldots, (t_{2k_2}, q_{2k_2})]),\\
  &\ldots,\\
  &(c_n, [(t_{n1}, q_{n1}), (t_{n2}, q_{n2}), \ldots, (t_{nk_n}, q_{nk_n})])]
\end{align*}

\noindent we let
$$
 \mathsf{valueData}(v) =
 \left\{
 \begin{array}{ll}
  m & \text{if $\Sigma_{i=1}^nk_i \leq 40000$}\\
  \errorX & \text{otherwise}
 \end{array}\right.
$$
\noindent
where

\begin{align*}
 m = \mathtt{Map}\,[
 &(\BB c_1, \mathtt{Map}\,[(\BB t_{11}, \II q_{11}), (\BB t_{12}, \II q_{12}), \ldots, (\BB t_{1k_1}, \II q_{1k_1})]),\\
 &(\BB c_2, \mathtt{Map}\,[(\BB t_{21}, \II q_{21}), (\BB t_{22}, \II q_{22}), \ldots, (\BB t_{2k_2}, \II q_{2k_2})]),\\
 &\ldots,\\
 &(\BB c_n, \mathtt{Map}\,[(\BB t_{n1}, \II q_{n1}), (\BB t_{n2}, \II q_{n2}), \ldots, (\BB t_{nk_n}, \II q_{nk_n})])].
\end{align*}


\note{Semantics of \texttt{unValueData}.}
\label{note:unValueData-semantics}
The \texttt{unValueData} function is the inverse of \texttt{valueData}.  It
takes a \texttt{data} object as input and returns a \texttt{value} object.  The
function requires its input to consist of a single outer \texttt{Map} whose
entries are pairs of (\texttt{B}, \texttt{Map}) pairs, with each inner map
containing a list of (\texttt{B}, \texttt{I}) pairs.  The keys in the outer map
and all of the inner maps must be at most 32 bytes long and they must be
strictly ascending according to the lexicographic ordering on bytestrings, and
all quantities (the \texttt{I} entries) must be nonzero and lie in the interval
$[-2^{127}, 2^{127}-1]$.  The outer map is allowed to be empty (representing the
empty value) but the inner maps must all be nonempty.

\medskip
\noindent
Formally, if $d$ is a data object then if $d$ is of the form

\begin{align*}
\mathtt{Map}\,[
 &(\BB c_1, \mathtt{Map}\,[(\BB t_{11}, \II q_{11}), (\BB t_{12}, \II q_{12}), \ldots, (\BB t_{1k_1}, \II q_{1k_1})]),\\
 &(\BB c_2, \mathtt{Map}\,[(\BB t_{21}, \II q_{21}), (\BB t_{22}, \II q_{22}), \ldots, (\BB t_{2k_2}, \II q_{2k_2})]),\\
 &\ldots,\\
 &(\BB c_n, \mathtt{Map}\,[(\BB t_{n1}, \II q_{n1}), (\BB t_{n2}, \II q_{n2}), \ldots, (\BB t_{nk_n}, \II q_{nk_n})])],
\end{align*}

\noindent then let 
\begin{align*}
  l = [
  &(c_1, [(t_{11}, q_{11}), (t_{12}, q_{12}), \ldots, (t_{1k_1}, q_{1k_1})]),\\
  &(c_2, [(t_{21}, q_{21}), (t_{22}, q_{22}), \ldots, (t_{2k_2}, q_{2k_2})]),\\
  &\ldots,\\
  &(c_n, [(t_{n1}, q_{n1}), (t_{n2}, q_{n2}), \ldots, (t_{nk_n}, q_{nk_n})])]
\end{align*}
\noindent and define 

$$
\mathsf{unValueData}(d) = \unflatten{l},
$$
remembering that $\unflatten{l} = \errorX$ if $l$ is not well formed.
If the \texttt{data} object $d$ is not of the expected form (involving
nested \texttt{Map}s) then also we let $\mathsf{unValueData}(d) = \errorX$.

\paragraph{Note.} Later we expect to extend the \texttt{data} type to include
a \texttt{Value} field containing a \texttt{value} object.  This will enable us
to introduce much simpler and faster versions of \texttt{valueData}
and \texttt{unValueData}.


\newcounter{notenumberG}
\renewcommand{\note}[1]{
  \bigskip
  \refstepcounter{notenumberG}
  \noindent\textbf{Note \thenotenumberG. #1}
}
\newpage
\subsection{Batch 7}
\label{sec:default-builtins-7}
A seventh batch of built-in functions is expected to be enabled at Protocol
Version 12.  They will be added here when their implementations and costing are
finalised.



